%%%%%%%%%%%%%%%%%%%%%%%%%%%%%%%%%%%%%%%%
%% PREAMBLE
%%%%%%%%%%%%%%%%%%%%%%%%%%%%%%%%%%%%%%%%
\documentclass[linenumbers,trackchanges,twocolumn]{aastex701}

\usepackage{amsmath}
\usepackage{multirow,booktabs,amssymb,xcolor,siunitx,makecell,rotating,xspace}

\newcommand{\abc}[3]{$#1/#2/#3$}
\newcommand{\vdag}{(v)^\dagger}
\newcommand\aastex{AAS\TeX}
\newcommand\latex{La\TeX}
\newcommand{\PRisma}{\texttt{PRisma}}
\graphicspath{{figures/}}

%%%%%%%%%%%%%%%%%%%%%%%%%%%%%%%%%%%%%%%%
%% MACROS
%%%%%%%%%%%%%%%%%%%%%%%%%%%%%%%%%%%%%%%%
\defcitealias{Zuluaga2015}{ZKSA2015}
\defcitealias{Kipping2014}{K2014}
\defcitealias{Seager2003}{SMO2003}

\newcommand{\PhotoRing}{\texttt{PhotoRing}}
\newcommand{\photofit}{PhotoFIT\xspace}

\newcommand{\parameters}{f_e, i_R, \theta_R, \tau,,p}

\newcommand{\observables}{\delta, T_{14}, T_{23}, b_{\mathrm{obs}}, (a/R_{\star})_{\rm obs}, \rho_{\star,\mathrm{obs}}}

\newcommand{\predictables}{\delta^\mathrm{geo}, T_{14}^\mathrm{geo}, T_{23}^\mathrm{geo}, b_{\mathrm{obs}}, (a/R_{\star})_{\rm obs}, \rho_{\star,\mathrm{obs}}}

\newcommand{\exoplanet}[2]{#1#2}

%%%%%%%%%%%%%%%%%%%%%%%%%%%%%%%%%%%%%%%%
%% COLORS
%%%%%%%%%%%%%%%%%%%%%%%%%%%%%%%%%%%%%%%%
\definecolor{jamm}{rgb}{0.018, 0.388, 0.878}
\newcommand{\JAAM}[1]{\textcolor{jamm}{#1}}%

\definecolor{numpaque}{rgb}{0.018, 0.388, 0.378} % 
\newcommand{\Numpaque}[1]{\textcolor{numpaque}{#1}}%

%\definecolor{attention}{rgb}{1,0,0} % 
\definecolor{attention}{rgb}{0,0,0} % 
\newcommand{\attention}[1]{\textcolor{attention}{#1}}%

%%%%%%%%%%%%%%%%%%%%%%%%%%%%%%%%%%%%%%%%
%% FRONTMATTER
%%%%%%%%%%%%%%%%%%%%%%%%%%%%%%%%%%%%%%%%
\begin{document}
% \title{Kepler-51b: A Super-Puff Ringed Exoplanet}
%\title{Is Kepler-51b, a Ringed Exoplanet?}
%\title{Asterodensity profile analysis of Kepler-51 planets: a study case for the detection of exoplanetary rings using PhotoRing effect}
%\title{Asterodensity profile analysis of the Kepler-51 planetary system:\\%
%a case study on the detection of exoplanetary rings via the \texttt{PhotoRing} effect}
%\title{Bayesian Retrieval of Planetary Rings in the Kepler-51 System via the PhotoRing Effect}
\title{Probing Exoplanetary Rings with Asterodensity Profiling: A PhotoRing Analysis of Kepler-51}
% \title{Stellar Densities as Probes of Planetary Rings: the \exoplanet{Kepler-51}{c}ase Study}

\author[0009-0000-5697-3416]{Sebasti\'an Numpaque}
\affiliation{SEAP/FACom, Instituto de F\'isica - FCEN, Universidad de Antioquia, Calle 70 No. 52-21, Medell\'in, Colombia.}
\email{david.rodriguez1@udea.edu.co} 

\author[0000-0002-6140-3116]{Jorge I. Zuluaga}
\affiliation{SEAP/FACom, Instituto de F\'isica - FCEN, Universidad de Antioquia, Calle 70 No. 52-21, Medell\'in, Colombia.}
\email{jorge.zuluaga@udea.edu.co}  

\author[0000-0002-4365-7366]{David Kipping}
\affiliation{Columbia University, 550 W 120th Street, New York NY 10027}
\email{dkipping@astro.columbia.edu} 

\author[0000-0003-0353-9741]{Jaime A. Alvarado-Montes} \altaffiliation{Macquarie University Research Fellow (MQRF)} 
\affiliation{Australian Astronomical Optics, Macquarie University, Balaclava Road, Sydney, NSW 2109, Australia.}
\affiliation{Astrophysics and Space Technologies Research Centre, Macquarie University, Balaclava Road, Sydney, NSW 2109, Australia.}
\email{jaime-andres.alvarado-montes@hdr.mq.edu.au}

\begin{abstract}
%Context
The Kepler-51 system hosts transiting super-puff planets whose extremely low bulk densities challenge conventional structural models, suggesting anomalously large inferred radii that might arise from unresolved planetary rings systems.
% Aim and Methods
We test this hypothesis by combining asterodensity profiling with a TTV-aware fit of transit light-curves of \textit{Kepler} long-cadence photometry. 
% Results
We find that the transit-inferred stellar densities for planets b and d systematically differ from independent isochrone-based estimates.
%---a discrepancy characteristic of the Before invoking rings, 
We discard the photoeccentric effect, demonstrating that the large eccentricities required to explain the density anomalies are robustly ruled out by current dynamical constraints. 
%Monte Carlo simulations show only a $\sim 3.3\%$ probability that Kepler-51b's anomaly is driven by unmodeled eccentricity, which is considerably lower than the $\sim 18.3\%$ probability for Kepler-51d. 
The density anomalies are also compatible with the \texttt{PhotoRing} (PR) effect caused by unmodeled rings. 
% Method
Using a geometric ringed-transit forward model and Bayesian retrieval, we identify configurations that reproduce the transit observables and account for the stellar density anomalies. Our solutions suggest compact, moderately inclined ring systems that increase the planets' true bulk densities by a factor of 3--4. 
% Conclusions
Rather than ultra-low-density bodies, these objects may be ``puff ringed planets,'' a scenario consistent with recent JWST observations of Kepler-51d. While the required impact parameters and true stellar densities show tension with independent constraints---which we attribute to Kepler's 30-minute cadence limiting the resolution of the contact geometry---this study presents the first fully developed application of the PR effect. We provide the complete open-source code to reproduce our results, support further follow-up of Kepler-51, and facilitate applications to other planetary systems.
\end{abstract}

\keywords{\uat{Exoplanet rings}{494} --- \uat{Transit photometry}{1709} --- \uat{Stellar properties}{1624} --- \uat{Monte Carlo methods}{2238}}

%%%%%%%%%%%%%%%%%%%%%%%%%%%%%%%%%%%%%%%%
%% CONTENT
%%%%%%%%%%%%%%%%%%%%%%%%%%%%%%%%%%%%%%%%
\section{Introduction} 

Several findings support the hypothesis that planetary rings are not exclusive to our Solar System but a natural outcome of possibly all planetary systems. In the Solar System, all four giant planets possess ring systems \citep{Charnoz2009} linked to tidal disruption of satellites, collisions among moonlets, or residual debris from planetary formation. More recently, the confirmation of rings around small bodies such as Chariklo \citep{Braga-Ribas2014}, Quaoar \citep{Morgado2023}, Chiron \citep{Ortiz2023}, and the dwarf planet Haumea \citep{Ortiz2017} suggests that ring formation can occur under a wide variety of physical conditions, raising more questions about how planetary rings are born and how they evolve. Still, exoplanetary rings (`exorings' for short) remain undiscovered among the thousands of confirmed extrasolar planets. Their detection represents a unique opportunity to explore and provide new constraints on formation and evolution of planetary systems. %JAAM

Early theoretical investigations demonstrated that Saturn-like ring systems could produce detectable signatures ($\leq 100$ ppm) in transit photometry \citep{Barnes2004}, especially during ingress and egress. However, although systematic searches for these signals in {\it Kepler} and {\it TESS} data have been conducted over the years \citep{Heising2015,Aizawa2018, Umetani2025RingsTESS}, no unambiguous detections of rings around transiting exoplanets have yet been confirmed \citep{Brown2001, Santos2015, Osborn2017, Ballesteros2018, LecavelierdesEtangs2017, Akinsanmi2020}. In practice, an obscuring structure around 1SWASP J140747b is presently considered the only confirmed exoring case \citep{Mamajek2012, Kenworthy2015}. However, given how young the system is, these ``rings'' are most likely composed of dust and gas and probably represent the remains of a circumplanetary disc, rather than debris from previous tidal interactions or moon disruptions. %JAAM

Exorings could also be detected via modern high-contrast imaging techniques \citep{Follette2023} or reflected-light photometry \citep{Arnold2004,Sucerquia2020,Zuluaga2022} and polarimetry \citep{Veenstra2025}. However, because these techniques are still limited, most studies have concentrated on identifying their distinctive imprints in stellar light curves, which change according to their orientation, radial extent, and optical depth (see eg. \citealt{Ohta2009,Zuluaga2015}).

The most salient effect of an exoring on a transit light-curve arises from the fact that the occulting area of a ringless planet can be significantly smaller than that of a ringed planet. If rings are not included in the transit model of the latter, the planetary radius will be overestimated (see e.g., \citealt{Zuluaga2015}). This would result in the bulk density being underestimated, which in turn could cause the planet to be incorrectly identified as an unusually low-density, “puffed” planet \citep{bakos2007hat}. In extreme cases, exorings could explain the anomalous densities of the so-called ``super-puffs'' or super puff planets  ($\rho<0.1$ g/cm$^3$, see eg. \citealt{Lee2016}), bodies with radii comparable to those of gas giants but only a few times the mass of Earth.

The planetary system Kepler-51 contains some of the most striking examples of so‑called super-puff planets. This compact, multi-planet system hosts at least three transiting planets—\exoplanet{Kepler-51}{b}, \exoplanet{Kepler-51}{c}, and \exoplanet{Kepler-51}{d}—with sizes between those of Neptune and Saturn, whose masses have been inferred from transit-timing variation (TTV) analyses (\citealt{Masuda2014}). With estimated masses of 5–6 $M_\oplus$, these planets are formally categorized as super-Earths. Their radii, however, are on the order of 6–10 $R_\oplus$, comparable to gas giants, which results in extremely low mean densities of $\rho < 0.06$ g cm$^{-3}$. As a consequence, they are currently the least dense—and thus the most extremely inflated—planets known. More recently, a fourth planet, Kepler-51 e, has been suggested \citep{Masuda2024}. Its presence is inferred indirectly, and its characteristics are consequently less well determined than those of the other planets in the system.

Several physical processes have been suggested to explain the extremely low densities and inflated radii of the Kepler-51 planets, including extended H-He envelopes, atmospheric escape and the possibility that they are still at a relatively early evolutionary stage (see, e.g., \citealt{Libby-Roberts2020}; \citealt{Wang2019}). Recent JWST observations of Kepler-51d, however, show, on the one hand, no significant photometric signature indicative of an equatorial bulge or a large opaque ring system \citep{Lammers2024}, and, on the other hand, a sloped transmission spectrum \citep{Libby-Roberts2025} that strongly supports the presence of a high-altitude haze layer, rendering the existence of a massive, opaque exoring highly unlikely. While an optically thin, very porous ring around Kepler-51d cannot be entirely excluded, prior theoretical work has already shown that standard rocky rings struggle to account for the extreme sizes of the Kepler-51 planets without assuming exceptionally porous material \citep{piro2020exploring}. In turn, the true nature of Kepler-51b remains unclear without further follow-up observations. Owing to this ongoing debate and their peculiar physical characteristics, the Kepler-51 planets serve as an excellent testing ground for the development and implementation of novel methods.

In \citet{Zuluaga2015} (hereafter \citetalias{Zuluaga2015}), the so‑called \PhotoRing\ (PR) effect was introduced, a technique for identifying planetary rings even when transit light curves lack prominent, well‑resolved morphological signatures. PR is part of the broader class of asterodensity profiling (AP) techniques (\citealt{Kipping2014}, hereafter \citetalias{Kipping2014}). In AP, the mean stellar density of the host star inferred from transit observables (the “observed density”)—namely the transit duration, depth, impact parameter, and orbital period— is compared with the bulk stellar density obtained from independent stellar characterization methods (the “true density”), such as asteroseismology or stellar isochrone-fitting. When the transit signal is influenced by unmodeled phenomena, such as stellar blending, orbital eccentricity, or the presence of planetary rings, the observed density will differ from the true density.

Despite the potential of the PR effect and the community’s positive reception of the concept, the method has not yet been applied to real planetary systems. With over a decade of accumulated observations, Kepler-51 now emerges as an optimal target to test the method and help unravel the baffling properties of its planets. To this end, we employ here PR to investigate whether the anomalous characteristics of the Kepler-51 planets can at least be partially explained due to the presence of planetary rings. We infer a stellar density from the Kepler transit light curves—among other transit observables—for each planet and compare it with an independent, isochrone-based estimate of the true stellar density \citep{Berger2023}, searching for discrepancies indicative of the PR effect. We then applied the geometric forward model for transiting ringed planets of \citetalias{Zuluaga2015} to predict the effective transit observables—namely, the transit depth $\delta$ and the durations $T_{14}$ and $T_{23}$—as well as the corresponding $\rho_{\star,\mathrm{obs}}$ for a specified ring configuration. %To determine which ring configurations are consistent with the Kepler data, we perform a Bayesian retrieval using {\em nested sampling}, which enables efficient exploration of the posterior distribution of ring parameters and assessment of the evidence for the ringed-planet model relative to a ringless model via the Bayes factor.

This paper is organized as follows: In \autoref{sec:geom_model}, we introduce the geometric forward model used to predict effective transit observables for ringed planets. \autoref{sec:asterodensity_profiling} outlines the asterodensity profiling framework and formally describes the \PhotoRing\ effect. In \autoref{sec:data_analysis}, we present the observational data analysis, including our light-curve fit of the Kepler-51 planets and the derivation of their anomalous transit-inferred stellar densities. \autoref{sec:bayes_inference} details the Bayesian inference methodology, explicitly describing the likelihood construction and the treatment of nuisance parameters. In \autoref{sec:ring_results}, we report the results of our retrieval setups, detailing the inferred ring geometries and evaluating the model through posterior predictive checks. Finally, \autoref{sec:discussion} provides a discussion of the physical implications of our findings and our concluding remarks.

\section{Geometric Modeling of Ringed Planet Transits} 
\label{sec:geom_model}

Our inference scheme relies on a deterministic, geometry-based forward model that connects a specified ring configuration to the associated effective transit observables. While we follow that relationship via the open-source \texttt{Python} package \texttt{exorings}\footnote{The original version is available at \url{https://github.com/facom/exorings}, and an updated release is maintained in this work's repository at \url{https://github.com/seap-udea/aPRe}.} \citepalias{Zuluaga2015}, this work represents an improved and expanded version of such analytical model on planetary ring transits.

The orbital period $P$, the impact parameter $b$, and the true stellar density $\rho_{\star,\mathrm{true}}$ are three key input parameters in the particular forward-model implementation adopted in this work. From these quantities, the geometric scaled semi-major axis is obtained via Kepler’s third law (\citealt{Seager2003}, hereafter \citetalias{Seager2003}):

\begin{equation}
\left(\frac{a}{R_\star}\right)_\mathrm{geo} = \left(\frac{G\,\rho_{\star,\mathrm{true}}}{3\pi}P^2\right)^{1/3}
\label{eq:aRstar_true}
\end{equation}
where $G$ is the gravitational constant. For a circular orbit ($e=0$), the orbital inclination is then given by 

\begin{equation}
\cos i_{\rm orb}=\frac{b}{(a/R_\star)_\mathrm{geo}}  
\label{eq:iorb}
\end{equation}

Two remarks are worth highlighting. First, by using the stellar density $\rho_{\star,\mathrm{true}}$ as an input parameter instead of the stellar mass $M_\star$ and radius $R_\star$, we avoid needing prior knowledge of those basic stellar properties for the analysis itself. Second, although we explicitly include $R_\star$ when referring to the scaled semi-major axis, all spatial scales are actually expressed in units of the stellar radius. As a result, the absolute value of $R_\star$ is not strictly required for the analysis. Still, to interpret the final inferred properties or express them in SI units, we must adopt values for both stellar parameters (or at least for one of them, given that the stellar density is taken to be known).

The ring system is modeled as a geometrically thin annulus extending from an inner radius $R_{\rm in}=f_i p$ to an outer radius $R_{\rm out}=f_e p$, where $p\equiv R_p/R_\star$ corresponds to the scaled planetary radius and $f_i, f_e$ are dimensionless radii in units of $p$. 

The orientation of the ring relative to the observer is described by two angles defined with respect to the plane of the sky. These angles determine the geometry of the ring’s sky-projected shape, which is modeled as a perfect ellipse with semimajor axis $A = f_e p$ and semiminor axis $B = A \cos i_R$. The orientation angles are $i_R$, the projected ring inclination (i.e., the dihedral angle between the ring plane and the sky plane, where $i_R = 90^\circ$ represents an edge-on view of the ring), and $\theta_R$, the projected \textit{tilt} angle (the angle between the projected major axis of the ring and the orbital direction, where $\theta_R = 90^\circ$ indicates that the projected major axis is perpendicular to the transit chord). For clarity, \autoref{tab:symbols} lists the definitions of all quantities used throughout this study.

\begin{table}[t]
\centering
\small
\renewcommand{\arraystretch}{1.2}
\caption{Summary of symbols and notation used throughout this work.}
\label{tab:symbols}
\begin{tabular}{ll}
\hline
\hline
Symbol & Description \\
\hline
\multicolumn{2}{l}{\textit{Nuisance parameters of the forward model}} \\
$P$ & Orbital period [hours] \\
$b$ & Transit impact parameter [--] \\
$\rho_{\star,\mathrm{true}}$ & Independent stellar density [g\,cm$^{-3}$] \\
\hline
\multicolumn{2}{l}{\textit{Ringed planet parameters of the forward model}} \\
$p$ & $R_p/R_\star$, scaled true planetary radius [--] \\
$f_i$ & Inner ring radius [$p$] \\
$f_e$ & Outer ring radius [$p$] \\
% $f_{e,\mathrm{max}}$ & Stability upper limit on outer ring radius [$p$] \\
$i_R$ & Projected ring inclination [deg] \\
$\theta_R$ & Projected ring tilt angle [deg] \\
$\tau$ & Ring normal optical depth [--] \\
$\beta$ & Ring blocking factor [--] \\
$A_\mathrm{eff}$ & Effective occulting area [$\pi R_\star^2$] \\
\hline
\multicolumn{2}{l}{\textit{Derived parameters of the forward model}} \\
$(a/R_{\star})_\mathrm{geo}$ & Geometric scaled semi-major axis[--] \\
$i_\mathrm{orb}$ & Orbital inclination [deg] \\
$\delta^\mathrm{geo}$ & Geometric transit depth [ppm] \\
$T_{14}^\mathrm{geo}$ & Geometric total transit duration [hours] \\
$T_{23}^\mathrm{geo}$ & Geometric full transit duration [hours] \\
\hline
\multicolumn{2}{l}{\textit{Transit (photometric) observables}} \\
$\delta$ & Transit depth [ppm] \\
$T_{14}$ & Total transit duration [hours] \\
$T_{23}$ & Full-transit duration [hours] \\
$p_\mathrm{obs}$ & Ringless-model scaled planetary radius [--] \\
\hline
\multicolumn{2}{l}{\it Asterodensity profiling derived properties} \\
$(a/R_{\star})_\mathrm{obs}$ & Scaled semi-major axis [--] \\
$\rho_{\star,\mathrm{obs}}$ & Transit-inferred stellar density [g\,cm$^{-3}$] \\
$\Psi$ & $\rho_{\star,\mathrm{obs}}/\rho_{\star,\mathrm{true}}$, density anomaly [--] \\
PR & $10\,\log_{10}\Psi$, logarithmic PR scale [--] \\
$\Theta_{ij}$ & $(\Psi_i/\Psi_j)^{2/3}$, Ratio of density anomalies [--]\\
&  (multiplanet asterodensity profiling, MAP) \\
\multicolumn{2}{l}{\textit{Bayesian inference}} \\
\hline
$\boldsymbol{\theta}$ & Ringed planet parameter vector  \\
$\boldsymbol{\eta}$ & Nuisance parameters  \\
$\mathbf{y}$ & Observable vector  \\
$\mathcal{L}$ & Likelihood function  \\
$\hat{p}(\mathbf{y})$ & KDE estimate of \photofit posterior  \\
\hline
\hline
\end{tabular}
\end{table}

Planetary rings are not fully opaque. To describe their opacity, we adopt a wavelength-independent normal optical depth $\tau$. This parameter is then used to derive a blocking factor $\beta = 1 - \exp(-\tau/\cos i_R)$ \citep{Barnes2004}, which incorporates both blocking and absorption effects and scales the ring’s contribution to the total occulting area according to its inclination.

Given these inputs, our forward model computes two classes of observables.  First, it evaluates the effective occulting area $A_{\rm eff}$ and the corresponding transit depth:

\begin{equation}
\delta^\mathrm{geo}=\frac{A_{\rm eff}}{\pi R_\star^2}.
\end{equation}

In practice, $A_{\rm eff}$ is written as the area of the planetary disk plus the effective projected area of the annulus. The latter is computed using closed-form expressions for the effective projected areas of the outer and inner ring edges as functions of $(f_i,f_e, i_R, \tau)$, including the edge-on limit where only a fraction of the annulus projects outside the planetary disk (Eq. 2 in \citetalias{Zuluaga2015}). Second, it determines the locations of the contact points along the transit chord. To illustrate the transit geometry with the exact definitions of the contact positions and transit durations, please refer to \autoref{fig:analytical_model}.

The planet-only contacts $x_{p,1\ldots4}$ are obtained from the intersection of the planetary disk with the stellar limb at impact parameter $b$. For the ring contacts $x_{R,1\ldots4}$, rather than adopting the approximate analytical intersection from \citetalias{Zuluaga2015} that breaks down at certain high inclinations, we employ a robust, exact-like analytical solution based on the Support Function of a convex body (see \autoref{sec:appendix_support_function}). The four ring-contact positions along the transit chord are given by:

\begin{eqnarray}
    x_{R,1} &\approx& -\sqrt{(1 + h_L)^2 - b^2} \nonumber \\
    x_{R,2} &\approx& -\sqrt{\max(0, [1 - h_L]^2 - b^2)} \label{eq:contact_x} \\
    x_{R,3} &\approx& +\sqrt{\max(0, [1 - h_R]^2 - b^2)} \nonumber \\
    x_{R,4} &\approx& +\sqrt{(1 + h_R)^2 - b^2}, \nonumber
\end{eqnarray}
where $h_L$ and $h_R$ are the effective directional radii of the projected ring ellipse at ingress and egress, respectively:
\begin{eqnarray}
    h_L^2 &=& A^2 (-x_0 \cos\theta_R + b \sin\theta_R)^2 \nonumber \\
          & & +\, B^2 (b \cos\theta_R + x_0 \sin\theta_R)^2 \\
    h_R^2 &=& A^2 (x_0 \cos\theta_R + b \sin\theta_R)^2 \nonumber \\
          & & +\, B^2 (b \cos\theta_R - x_0 \sin\theta_R)^2 \nonumber
\end{eqnarray}
(see equations \ref{eq:h_L} and \ref{eq:h_R} in the Appendix). These formulae are key to the methodological implementation used here: as they are fully analytical, we can perform computationally intensive numerical procedures without incurring prohibitively long computation times.

The final contact positions are then taken as $x_i=\min[x_{p,i},x_{R,i}]$, and converted into the ringed transit durations via
\begin{equation}
T_{14}^\mathrm{geo} = \frac{P}{2\pi}\arcsin\!\left[\frac{x_{4}-x_{1}}{(a/R_\star)_\mathrm{geo}\sin i_{\rm orb}}\right]
\label{eq:T14_geo}
\end{equation}
\begin{equation}
T_{23}^\mathrm{geo} = \frac{P}{2\pi}\arcsin\!\left[\frac{x_{3}-x_{2}}{(a/R_\star)_\mathrm{geo}\sin i_{\rm orb}}\right].
\label{eq:T23_geo}
\end{equation}

This mathematical procedure specifies the forward mapping from the ringed-planet parameters $(p, f_i, f_e, i_R, \theta_R, \tau)$ to the predicted effective transit observables $(\delta^\mathrm{geo}, T_{14}^\mathrm{geo}, T_{23}^\mathrm{geo})$, under the assumption that the stellar and orbital input parameters $P$, $b$, and $\rho_{\star,\mathrm{true}}$ are independently provided.

\section{Asterodensity Profiling}
\label{sec:asterodensity_profiling}

Assuming a spherical, opaque planet with a mass $M_\mathrm{ p}\ll M_\star$, in a circular orbit and transiting a spherical, unblended star of radius $R_\star$ with negligible limb darkening at the maximum transit depth position, \citetalias{Seager2003} showed a mean stellar density, $\rho_{\star,\mathrm{obs}}$, can be inferred from the transit observables by inverting~\autoref{eq:aRstar_true}:

\begin{equation}
\rho_{\star,\mathrm{obs}} = \frac{3\pi}{GP^2} \left(\frac{a}{R_\star}\right)^3_{\mathrm{obs}}.
\label{eq:rho_obs}
\end{equation}

The observed scaled semi-major axis $(a/R_\star)_\mathrm{obs}$ should not be mistaken for $(a/R_\star)_\mathrm{geo}$ that is computed in the forward model from the orbital period and the true stellar density (\autoref{eq:aRstar_true}). \citetalias{Seager2003} demonstrated that $(a/R_\star)_\mathrm{obs}$ can be expressed in terms of the total $T_{14}$ and full-transit duration $T_{23}$, via their Equation 8\footnote{We distinguish here between the {\em measured} transit parameters $T_{14}$ and $T_{23}$ and the {\em calculated} values $T_{14}^\mathrm{geo}$ and $T_{23}^\mathrm{geo}$ produced by the forward model}:

\begin{equation}
\left(\frac{a}{R_\star}\right)^2_{\mathrm{obs}} = \frac{F_{(+)}^2 - b_{\mathrm{obs}}^2 \left[1 - \sin^2\left(\pi T_{14}/P\right)\right]}{\sin^2\left(\pi T_{14}/P\right)},
\label{eq:aRstar}
\end{equation}
where $F_{(\pm)}\equiv (1 \pm \sqrt{\delta})$ is an auxiliary variable. 

This formula assumes non-grazing transits ($1<b<1-p$), a planet that is much smaller than its host star ($p \ll 1$), and transit durations that are much shorter than the orbital period ($T_{14}/P \ll 1$). These conditions are typically met in most transiting exoplanet systems \citepalias{Seager2003, Kipping2014}, including the Kepler-51 planets.

Under the same assumptions, the observed impact parameter $b_{\mathrm{obs}}$ is obtained from the transit shape via Equation 9 in \citetalias{Seager2003}:

\begin{equation}
b_{\mathrm{obs}}^2 = \frac{F_{(-)}^2 - \sin^2\left(\pi T_{23}/P\right) / \sin^2\left(\pi T_{14}/P\right)F_{(+)}^2}{1 - \sin^2\left(\pi T_{23}/P\right) / \sin^2\left(\pi T_{14}/P\right)}.
\label{eq:b_obs}
\end{equation}

It is worth noting that equations (\ref{eq:aRstar}) and (\ref{eq:b_obs}), and thus \autoref{eq:rho_obs}, are equally valid when using the geometrically computed durations $T_{14}^\mathrm{geo}$ and $T_{23}^\mathrm{geo}$ from the forward model, obtained via \autoref{eq:T14_geo} and \autoref{eq:T23_geo}. The specific context in which the formulae are employed determines which set of quantities should be used. For example, when deriving the observed semi-major axis in the photometric analysis (see \autoref{sec:derived_observables}), the time and period values are those inferred from the photometric data, namely $\delta$, $T_{14}$, and $T_{23}$. In contrast, during the inference stage we employ instead the geometric quantities $\delta^\mathrm{geo}$, $T_{14}^\mathrm{geo}$, and $T_{23}^\mathrm{geo}$ (see \autoref{sec:bayes_inference}).

The concept of asterodensity profiling \citepalias{Kipping2014} is based on the comparison between the photometrically inferred stellar density, $\rho_{\star,\mathrm{obs}}$, and an independent estimate of the true stellar density, $\rho_{\star,\mathrm{true}}$, obtained through methods such as spectroscopic analysis, asteroseismology, or stellar evolution modeling. If all assumptions of the simple transit model are satisfied, both density estimates should agree within observational uncertainties.

Asterodensity profiling also explores various astrophysical mechanisms for diagnosing violations of the idealized transit assumptions through statistically significant deviations between $\rho_{\star,\mathrm{obs}}$ and $\rho_{\star,\mathrm{true}}$ One well-known example is the photoeccentric effect, which arises when a planet on an eccentric orbit is incorrectly modeled as circular \citepalias{Kipping2014}. Other mechanisms include contamination by blended companions (photoblend), unocculted stellar activity (photospot), and additional transit-shape distortions such as those induced by planetary rings (photoring, \citetalias{Zuluaga2015}).


\subsection{The \PhotoRing\ Effect}
\label{sec:photoring_effect}

Planetary rings introduce geometric effects that are not captured by the standard spherical-planet transit model. As shown by \citetalias{Zuluaga2015}, the projected area of a ring system can substantially increase the observed transit depth $\delta_{\mathrm{obs}}$. When rings are not included in the transit model, this larger depth $\delta_{\mathrm{obs}}$ is interpreted as being produced by a larger planetary disk, which leads to an overestimation of the planetary radius $R_p$.

Rings also modify the contact times of a transit. Because the ring system extends beyond the planetary limb, the first contact occurs earlier and the fourth contact occurs later than in the case of a ringless planet. As a result, the total transit duration $T_{14}$ increases, while the full-transit duration $T_{23}$ becomes shorter. Since these three quantities constitute the fundamental transit observables, these geometric distortions propagate through the transit fitting procedure and bias the inferred values of $(a/R_\star)_{\mathrm{obs}}$ and $b_{\mathrm{obs}}$ (i.e., equations \ref{eq:aRstar} and \ref{eq:b_obs}, respectively). Consequently, the stellar density derived from the transit, $\rho_{\star,\mathrm{obs}}$, can be significantly under or overestimated. The resulting systematic offset in the stellar density inferred from transits due to an unaccounted-for ring system was first characterized by \citetalias{Zuluaga2015} and is known as the \PhotoRing\ (PR) effect.

By introducing the stellar anomaly ratio, following \citet{kipping2012novel}, as

\begin{equation}
\Psi \equiv \left(\frac{\rho_{\star,\mathrm{obs}}}{\rho_{\star,\mathrm{true}}}\right),
\label{eq:PR_Kipping}
\end{equation}
the PR effect is expressed using a logarithmic scale that employs the same notation:

\begin{equation}
\mathrm{PR} \equiv 10\,\log_{10}\!\Psi.
\label{eq:PR_def}
\end{equation}

In \autoref{fig:photoring_contour}, we present contours of the PR-effect metric across a wide range of ring orientations. We adopt planetary and stellar parameters akin to those of \exoplanet{Kepler-51}{b}, along with tentative properties for the ring system. As shown, PR effect typically causes the stellar density to be underestimated; that is, $\rho_{\star,\mathrm{obs}} < \rho_{\star,\mathrm{true}}$ or $\mathrm{PR}<0$, with only a relatively small domain where $\rho_{\star,\mathrm{obs}} > \rho_{\star,\mathrm{true}}$. Interestingly, for systems hosting ringed planets—with similar properties to those assumed here—the inferred stellar mean density can be underestimated by as much as a factor of $\approx 2.5$ ($\mathrm{PR}\approx -4$) or overestimated by up to a factor of $\approx 1.6$ ($\mathrm{PR}\approx +2$).

\begin{figure*}[t]
\centering
\includegraphics[width=0.7\textwidth]{figures/PRContour-K51b-exorings_updated.png}
\caption{Heat map and contour lines of the \PhotoRing\ (PR) logarithmic scale as a function of the effective ring orientation parameters $(i_{\rm R},\theta_{\rm R})$. PR is evaluated for an analogue of \exoplanet{Kepler-51}{b} with a bulk density equal to that of Earth. The thick black contour marks the zero level, $PR=0$, which identifies configurations that exhibit a density degeneracy (i.e., the inferred stellar density is equal to the true density despite the presence of rings). The highlighted contour correspond to the nominal observed anomaly for the actual planet \exoplanet{Kepler-51}{b} according to our \photofit fits (see text for details).}
\label{fig:photoring_contour}
\end{figure*}

% In systems where blended companions can be ruled out, for instance via high-resolution imaging, these deviations can be separated from other asterodensity profiling (AP) effects. Specifically, the PhotoEccentric effect (PE) generally generates a positive AP deviation because of the bias introduced by the geometric transit probability \citepalias{Kipping2014}. Consequently, a negative offset between $\rho_{\star,\mathrm{obs}}$ and $\rho_{\star,\mathrm{true}}$ may indicate the presence of rings rather than orbital eccentricity or blending.

\section{Observational Data Analysis}
\label{sec:data_analysis}

The observational and data foundation of this work rests on posterior distributions of the asterodensity profiling observables ($\delta, \rho_{\star,\mathrm{obs}},  b_{\rm obs}$) for \exoplanet{Kepler-51}{b} and \exoplanet{Kepler-51}{d}, the innermost and outermost planets of Kepler-51, obtained from a dedicated light-curve fit of the available \textit{Kepler} photometry (see \autoref{sec:data_analysis_ttv}). These posteriors are not derived from a direct ring-transit model, they rather represent the inferences made under the standard assumption of a spherical, ringless planet, thus encoding the effective transit observables that must result from any ringed-planet configuration.

In this section, we describe the origin of these posteriors and how they are incorporated into our ring-retrieval framework. Although more recent observations of the system have been obtained with JWST, HST, and TESS facilities \citep{Masuda2024, Lammers2024, Libby-Roberts2025}, the present analysis exclusively relies on constraints derived from the original \textit{Kepler} photometry, without including the aforementioned new datasets.

\subsection{Transit and \photofit Modeling}
\label{sec:data_analysis_ttv}

For a robust inference of the transit parameters in the presence of TTVs, without resorting to a full dynamical analysis, we employed a so-called {\em photometric fit with independent transit times} (hereafter \photofit). Enforcing a strict linear ephemeris on systems with significant TTVs can artificially smear the transit shape. In our \photofit approach, the central transit time ($T_c$) for each epoch is treated as an independent free parameter during the light curve fitting process. This decouples the transit's photometric shape from the underlying orbital dynamics, allowing us to accurately measure the geometric transit parameters while also extracting the individual transit epochs.

For \exoplanet{Kepler-51}{b} and \exoplanet{Kepler-51}{d}, the asterodensity profiling posterior distributions were sampled using the nested sampling algorithm \texttt{MultiNest} \citep{Feroz2009, Feroz2019} applied to the \textit{Kepler} photometric time series\footnote{We exclude the recently discovered fourth planet, Kepler-51 e, and \exoplanet{Kepler-51}{c} due to its grazing transit geometry \citep{Masuda2024}}. The analysis fits each transit simultaneously with our \photofit model, extracting individual mid-transit times for each observed epoch while allowing to estimate the AP parameters from the full ensemble of transits.

\attention{For \exoplanet{Kepler-51}{d}, the small number of observed transits and the corresponding parameter space allowed a single nested-sampling run, yielding one posterior distribution per parameter. For \exoplanet{Kepler-51}{b}, however, the large number of observed epochs translates into a high-dimensional parameter space---since each epoch contributes an individual mid-transit time as a free parameter---making a single nested-sampling run computationally intractable. To handle this high dimensionality, the full photometric time series was divided into three independent temporal segments, each covering a distinct subset of observed transits. A separate run was performed on each segment, producing its corresponding posterior distributions for AP parameters. We then  combine the resulting distributions into a single posterior by running a short Markov Chain Monte Carlo sampler that simultaneously draws from the three posterior distributions, treating them as statistically independent constraints on the transit parameters. The resulting chain defines the final posterior of \exoplanet{Kepler-51}{b} used throughout this work.}

\subsection{Derived Observables}
\label{sec:derived_observables}

For each planet, the main outputs of the \photofit are joint posterior samples of the AP parameters $(\rho_{\star,\mathrm{obs}}, b_{\mathrm{obs}}, \delta)$, along with the orbital period $P$ and the limb-darkening coefficients. Assuming the planet has no rings, these parameters capture all of the transit-geometry information encoded in the \textit{Kepler} light curves. The secondary transit quantities $(a/R_\star)_{\mathrm{obs}}, T_{14}$, and $T_{23}$, which are needed for our AP analysis, are then obtained by deterministically inverting the relations given in \citetalias{Kipping2014}. Their posterior distributions are derived directly from the same posterior samples, ensuring that the full covariance structure inferred from the photometry is preserved. In \autoref{fig:observables_full}, we show the marginal posterior distributions for the full set of observables $(\observables)$ for both \exoplanet{Kepler-51}{b} and \exoplanet{Kepler-51}{d}. Our results (shaded curves) are compared with those obtained by other authors (see next section).

\begin{figure*}[t]
\centering
% Replace with actual figure path when available
\includegraphics[width=\textwidth]{k51_observables_posteriors.png}
\caption{Marginal posterior distributions of the observables for \exoplanet{Kepler-51}{b} (top row) and \exoplanet{Kepler-51}{d} (bottom row) obtained from the \photofit carried out in this work are shown as filled distributions. Vertical dashed lines and the accompanying labels indicate the median and the 68\% credible interval for each marginal posterior. For comparison, literature estimates are depicted as normal distributions (unfilled curves) constructed from the published point values and their reported uncertainties; in the case of \citet{Masuda2024}, a truncated normal distribution is used for the impact parameter $b$ to respect its physical limits. Only those observables provided by each reference are plotted (for example, none of the alternative sources report values for the flat-bottom duration $T_{23}$).}
\label{fig:observables_full}
\end{figure*}

As detailed in \autoref{sec:geom_model} and \autoref{sec:asterodensity_profiling}, the geometrical forward model provides, for a specified set of ring parameters $(p, f_i, f_e, i_R, \theta_R, \tau)$ and suitable choices of orbital and stellar properties $(P, b, \rho_\mathrm{\star,true})$, a prediction for the corresponding observable vector. The Bayesian inference on the ring parameters is subsequently carried out by confronting this model prediction with the distribution derived from the light curves. This procedure exactly constitutes the foundation of the PR analysis. 

% -------------------------------------------------------------------
\subsection{Literature Comparison}
\label{sec:data_analysis_comparison}
% -------------------------------------------------------------------

The Kepler-51 system has been characterized by several independent photometric analyses spanning a decade: the confirmed three-planet architecture by \citet{Masuda2014}, the atmospheric characterization of \citet{Libby-Roberts2020} and \citet{Libby-Roberts2025}, and the refined dynamically re-analysis and discovery of a fourth non-transiting planet, \exoplanet{Kepler-51}{e} of \citet{Masuda2024}. Since our goal is to obtain posterior distributions of the photometric transit observables that enter the asterodensity profiling and ring inference described in \autoref{sec:asterodensity_profiling}--\autoref{sec:bayes_inference}, the comparison between our \photofit-derived posteriors and the literature point estimates therefore serves two purposes: (i) to validate that our analysis recovers transit shapes consistent with independent measurements, and (ii) to identify and explain genuine discrepancies that may arise from differences in data selection, dynamical assumptions, or statistical frameworks. In \autoref{tab:observables_full} we summarize the values of the transit observables derived from our \photofit, along with the corresponding literature values. 

A visual comparison of the values reported in the literature with those derived from our photometric analysis is presented in \autoref{fig:observables_full}. Because the posterior distributions of the observables reported by other authors are not available, we assume in all cases that they follow Gaussian distributions (shown as empty curves). While our own posteriors demonstrate that observable distributions can deviate substantially from normality, this assumption is made solely for comparative purposes and the exact shapes of the literature distributions are not used in our inference pipeline.

% Auto-generated by k51_observables_plotting.ipynb — do not edit by hand
% Requires \usepackage{multirow} in the preamble
\begin{table*}[t]
\footnotesize
\centering
\renewcommand{\arraystretch}{1.2}
\setlength{\tabcolsep}{15pt} % Default is 6pt
\caption{Compendium of transit observables for \exoplanet{Kepler-51}{b} and d: this work (\photofit posterior) and literature values}
% , and the \citet{Berger2023} isochrone-based stellar density estimation. The values reported for $\rho_{\star, \rm obs}$ are derived through \autoref{eq:rho_obs}.}
\label{tab:observables_full}
\begin{tabular}{llcc}
\hline\hline
Observable & Source & \exoplanet{Kepler-51}{b} & \exoplanet{Kepler-51}{d} \\
\hline
\multirow{5}{*}{$p$ [$R_{\star}$ ($R_{\oplus}$)]} & This work (\photofit) & $0.07055_{-0.00030}^{+0.00026}\,(6.687_{-0.028}^{+0.024})$ & $0.09531_{-0.00078}^{+0.00082}\,(9.034_{-0.074}^{+0.078})$ \\
 & \citet{Masuda2024} & $0.07225 \pm 0.0003 \,(6.83 \pm 0.13)$ & $0.09857 \pm 0.00037\,(9.32 \pm 0.18)$ \\
 & \citet{Masuda2014} & $0.07414_{-0.00061}^{0.00059}\,  (7.1\pm 0.3)$ & $0.10141_{-0.00085}^{0.00084} \, (9.7 \pm 0.5)$ \\
 & \citet{Libby-Roberts2020} & $0.0689\pm0.0013\,(6.89\pm0.14)$ & $0.0976\pm0.0009\,(9.46\pm0.16)$ \\
 & \citet{Libby-Roberts2025} & -- & $0.09682 \pm 0.00051\, (9.19 \pm 0.05)$ \\
\hline
\multirow{5}{*}{$\delta$ [ppm]} & This work (\photofit) & $4977_{-42.1}^{+36.1}$ & $9084_{-149}^{+156}$ \\
 & \citet{Masuda2024} & $5220 \pm 40$ & $9710 \pm 70$ \\
 & \citet{Masuda2014} & $5500 \pm 90$ & $10280 \pm 170$ \\
 & \citet{Libby-Roberts2020} & $4750 \pm 180$ & $9530 \pm 180$ \\
 & \citet{Libby-Roberts2025} & -- & $9370 \pm 100$ \\
\hline
\multirow{2}{*}{$T_{14}$ [h]} & This work (\photofit) & $5.708_{-0.0176}^{+0.0238}$ & $8.458_{-0.036}^{+0.0377}$ \\
 & \citet{Masuda2024} & $5.754 \pm 0.02$ & $8.402 \pm 0.029$ \\
\hline
\multirow{1}{*}{$T_{23}$ [h]} & This work (\photofit) & $4.939_{-0.019}^{+0.0243}$ & $6.88_{-0.0752}^{+0.0686}$ \\
% \hline
% \multirow{5}{*}{$\rho_{\star,\rm obs}$ [g cm$^{-3}$]} & This work (\photofit) & $2.441_{-0.0732}^{+0.0529}$ & $2.171_{-0.147}^{+0.158}$ \\
%  & \citet{Masuda2014} (derived) & $2.158 \pm 0.142$ & $2.165 \pm 0.143$ \\
%  & \citet{Libby-Roberts2020} (derived) & $1.974 \pm 0.288$ & $2.175 \pm 0.199$ \\
%  & \citet{Libby-Roberts2025} (derived) & -- & $2.139 \pm 0.019$ \\
%  & \citet{Berger2023} ($\rho_{\star,\rm true}$) & $1.896_{-0.246}^{+0.218}$ & $1.896_{-0.246}^{+0.218}$ \\
\hline
\multirow{4}{*}{$b_{\rm obs}$} & This work (\photofit) & $0.1443_{-0.0579}^{+0.0551}$ & $0.2696_{-0.116}^{+0.0757}$ \\
 & \citet{Masuda2024} & $0.074 \pm 0.072$ & $0.003 \pm 0.095$ \\
 & \citet{Masuda2014} & $0.251 \pm 0.106$ & $0.25 \pm 0.108$ \\
 & \citet{Libby-Roberts2020} & $0.22 \pm 0.16$ & $0.19 \pm 0.145$ \\
\hline
\multirow{4}{*}{$a/R_{\star}$} & This work (\photofit) & $64.08_{-0.647}^{+0.46}$ & $124.8_{-2.88}^{+2.96}$ \\
 & \citet{Masuda2024} & $61.5 \pm 1.35$ & $124.7 \pm 2.75$ \\
 & \citet{Libby-Roberts2020} & $59.7 \pm 2.9$ & $124.9 \pm 3.8$ \\
 & \citet{Libby-Roberts2025} & -- & $124.2 \pm 0.37$ \\
%  & \citet{Berger2023} (derived) & $58.9_{-2.66}^{+2.18}$ & $119.3_{-5.4}^{+4.41}$ \\
\hline
\end{tabular}
\end{table*}

The transit depth $\delta$ is the most straightforward quantity to determine from the data. \autoref{tab:observables_full} and \autoref{fig:observables_full} reveal a pronounced scatter among published values for both planets that is larger than the quoted uncertainties of the individual measurements. The most natural explanation for this spread is \emph{unocculted stellar activity}, rather than any systematic issue in the light-curve fitting. This phenomenon was explicitly detected in the follow-up observations of \exoplanet{Kepler-51}{d} using NIRspec, as documented by \citet{Lammers2024}. Kepler-51 is a young ($\sim$500\,Myr), magnetically active star whose evolving starspot distribution alters the apparent out-of-transit flux level, artificially increasing the measured transit depth if the spot contribution is not explicitly modeled. Since the spot covering fraction changes over the \textit{Kepler} observing window, and our analysis does not apply a dedicated correction for stellar spots, the derived depths should be interpreted as an average over the star’s activity state across the different epochs. This explains their agreement with the full span of literature values and supports the conclusion that the observed dispersion is predominantly caused by time-varying stellar activity, rather than by differences in analysis methodology.

On the other hand, the impact parameter $b_{\mathrm{obs}}$ is among the least constrained transit observables reported in the literature, with published estimates spanning a wide range for both planets. This behavior arises primarily from the limited temporal resolution of the \textit{Kepler} long-cadence data ($\sim$30 min), which is comparable to the ingress/egress durations of the planets and therefore averages over and smooths the slopes of these regions that contain most of the information used to constrain the transit chord geometry and, therefore, the impact parameter. This loss of geometric information introduces a strong degeneracy in the determination of $b_{\rm obs}$ that can only be effectively broken with higher-cadence observations, as demonstrated by the JWST analysis of \citet{Masuda2024}, which resolves a nearly central value for \exoplanet{Kepler-51}{d} ($b \sim 0.003$) and also show that, even with high-quality data, the solution for the impact parameter of \exoplanet{Kepler-51}{b} remains bimodal, implying the existence of two distinct transit geometries that fit the observations equally well. Our posterior for $b_{\rm obs}$, obtained from a joint fit of all available epochs, is consistent with both the low-$b$ and high-$b$ solutions reported and shows no significant tension with individual measurements, thereby reflecting the degeneracy imposed by the temporal resolution of the available photometry instead of a physical disagreement.

The total transit duration $T_{14}$ and the flat-bottom duration $T_{23}$ are not frequently tabulated in the literature, because most works describe the transit shape directly using $(a/R_\star)_{\mathrm{obs}}$ and $b_{\mathrm{obs}}$ rather than the contact-time durations themselves. In the one case where $T_{14}$ has been published by \citet{Masuda2024}, our posterior values for both planets are consistent, differing by at most $\lesssim4$ minutes, which lies comfortably within the quoted uncertainties. In contrast, the scaled semi-major axis $(a/R_\star)_{\mathrm{obs}}$ is the most consistently determined observable across all datasets and analyses: our measurements for both planets agree with those of \citet{Masuda2024}, \citet{Libby-Roberts2020}, and the independent JWST-based determination of \citet{Libby-Roberts2025} to within $1\sigma$. This agreement shows that the orbital scale of the system is effectively unaffected by the cadence, wavelength coverage, or modeling assumptions described above, and that $(a/R_\star)_{\mathrm{obs}}$ is set almost entirely by the precisely measured orbital period and total transit duration.

\subsection{Anomalous Stellar Densities}
\label{sec:anomalous_stellar_densities}

A key consequence of the asterodensity profiling framework is that, within the standard transit model—assuming unblended, spot-free light curves and spherical planets on circular orbits—all transiting planets in the same system must yield an identical value of $\rho_{\star,\mathrm{obs}}$. This conclusion follows directly from \autoref{eq:rho_obs}, since the inferred stellar density is computed from the measured orbital scale $(a/R_\star)_{\mathrm{obs}}$, which in turn depends only on the transit light-curve morphology and the orbital period. Consequently, any discrepancy in $\rho_{\star,\mathrm{obs}}$ among multiple planets orbiting the same star indicates that the standard transit model is failing to capture some aspect of the actual transit geometry for at least one of the planets, or that the transit light curves are being systematically modified by a physical effect not accounted for in the model.. %RF

\autoref{fig:rhoobs_comparison} presents the marginal posterior distributions of $\rho_{\star,\mathrm{obs}}$ derived from our photometric analysis for both planets, together with the corresponding distribution of $\rho_{\star,\mathrm{true}}$ from the independent isochrone-based estimates of \citet{Berger2023} (shown as filled curves). For comparison, we also display the values reported by \citet{Libby-Roberts2020} and \citet{Masuda2024}, modeling their posteriors as normal distributions based on their quoted means and uncertainties. For ease of comparison, the mean values and associated 1-$\sigma$ or 68\% uncertainties from our work and from the literature are compiled in \autoref{tab:rho_star_compendium}. %RF

\begin{figure*}[t]
\centering
% Replace with actual figure path when available
\includegraphics[width=\textwidth]{rho_obs_comparison.png}
\caption{Shown are the marginal posterior distributions of the transit-derived stellar density, $\rho_{\star,\mathrm{obs}}$, for \exoplanet{Kepler-51}{b} (blue) and d (orange), obtained from the \photofit presented here. The vertical dashed lines and accompanying labels mark the median and the 68$\%$ credible interval for each posterior. The discontinuous curves correspond to Gaussian distributions constructed from literature point estimates and their reported uncertainties from \citet{Libby-Roberts2020} (purple) and \citet{Masuda2024} (green). Also plotted are the independent isochrone-based posterior samples for the true stellar density, $\rho_{\star,\mathrm{true}}$, from \citet{Berger2023} (filled pink).} %RF
\label{fig:rhoobs_comparison}
\end{figure*}

\begin{table*}[t]
\centering
\renewcommand{\arraystretch}{1.2}
\caption{Estimates of the stellar density for the Kepler-51 system. The values are grouped according to whether they are obtained from transit observables (including photodynamical modeling) or from independent stellar characterization. The transit-based estimates reported by \citet{Masuda2024} were derived under the assumption of the previously known three-planet configuration, whereas their isochrone-based stellar characterization was carried out using the current four-planet model. For reference, we also list an estimate of the logarithmic PR scale (\autoref{eq:PR_def}), computed by adopting a true reference stellar density of $2$ g cm$^{-3}$, which is close to the mean of the independent estimates. The uncertainties on PR were computed via straightforward linear error propagation. The PR values are reported for context only and are not used in the inference analysis. \label{tab:rho_star_compendium}}
\begin{tabular}{lllcc}
\hline\hline
Category & Source & Method & $\rho_{\star}$ [g cm$^{-3}$] & PR (ref.) \\
\hline

\multirow{6}{*}{Transit-derived}
& This work (\exoplanet{Kepler-51}{b})
& \photofit posterior
& $2.441_{-0.073}^{+0.053}$
& $0.87\pm 0.12$
\\

& This work (\exoplanet{Kepler-51}{d})
& \photofit posterior
& $2.171_{-0.147}^{+0.158}$ 
& $0.36\pm 0.31$
\\

& \citet{Masuda2014}
& Transit light-curve fit
& $2.16_{-0.13}^{+0.15}$ 
& $0.33\pm 0.28$
\\

& \citet{Masuda2024}
& Transit light-curve fit
& $2.42 \pm 0.06$ 
& $0.83\pm 0.11$
\\

& \citet{Masuda2024}
& N-body Photodynamical fit
& $\sim 2.56_{-0.34686}^{+0.423}$ 
& $1.07\pm 0.65$
\\

\hline

\multirow{4}{*}{Independent}

& \citet{Libby-Roberts2020}
& Isochrone fitting (CKS+Gaia DR2)
& $2.03 \pm 0.08$ 
& --\\

& \citet{Masuda2024}
& Isochrone fitting (MIST+Gaia DR3)
& $2.08 \pm 0.08$ 
& --\\

& \citet{Berger2023}
& Isochrone fitting (Gaia DR3)
& $1.896_{-0.246}^{+0.218}$ 
& --\\
\hline
\end{tabular}
\end{table*}

Several mechanisms could explain the discrepancy between the stellar density inferred from photometry and that obtained from independent stellar-astrophysics methods: 

\begin{itemize}
    \item[(i)] an underlying non-circular orbit, commonly known as the photoeccentric effect \citep{Dawson2012,kipping2012novel,Kipping2014}. If a planet transits near periastron, its higher orbital velocity mimics a closer orbit, leading to an artificially inflated inferred stellar density.
    
    \item[(ii)] the omission of an explicit spot correction in the photometric analysis, which affects the measured $(a/R_\star)_{\mathrm{obs}}$ (\autoref{eq:aRstar}) and, in turn, $\rho_{\star,\mathrm{obs}}$. In our study, this effect could partially explain the higher stellar densities derived for both planets, as outlined in \autoref{sec:data_analysis_comparison}. %RF
    
    \item[(iii)] Unmodeled dynamical interactions among the planets. Kepler-51 exhibits large-amplitude transit timing variations (TTVs) that can distort the effective transit light-curve shape if not properly accounted for. In our analysis, we substantially mitigate this bias by fitting a separate mid-transit time for each individual epoch. More generally, an incomplete characterization of the system’s architecture can also bias stellar-density estimates. This is demonstrated by the three-planet configuration assumed by \citet{Masuda2014}, as well as by both the transit-fitting and photodynamical analyses of \citet{Masuda2024} (\autoref{tab:rho_star_compendium}). That truncated dynamical model incorrectly predicted the \(JWST\)-measured mid-transit time of Kepler-51d by \(\sim 120\) minutes. Consequently, stellar-density inferences that rely on this assumption—including those presented here, which likewise neglect perturbations from the fourth planet—may also be affected by systematic biases. %RF
    
    \item[(iv)] The existence of an extra geometric feature—such as a ring system—that changes the observed transit configuration by enlarging the effective occulting area and altering the transit contact times. %RF
\end{itemize}

\subsection{Ruling out the Photoeccentric Effect}
\label{sec:photoeccentric_effect}

The positive density anomaly observed in Kepler-51b and d ($\Psi \equiv \rho_{\star,\mathrm{obs}}/\rho_{\star,\mathrm{true}} > 1$) could in principle be induced by an unmodeled orbital eccentricity if the planets transit near periastron. This is known as the photoeccentric effect \citep{Dawson2012}. To verify whether this mechanism can account for the observed $\rho_{\star,\mathrm{obs}}$, we calculate the minimum eccentricity $e_{\mathrm{min}}$ required by the \photofit and compare it against the dynamical constraints imposed by TTVs \citep{Kipping2014}. 

For a single planet, the minimum eccentricity capable of producing an observed density ratio $\Psi > 1$ (which occurs when the argument of periastron $\omega = 90^\circ$) is given by:
\begin{equation}
    e_{\mathrm{min}} = \frac{\Psi^{2/3} - 1}{\Psi^{2/3} + 1}.
\end{equation}

By evaluating this expression over our posterior samples of $\Psi$ for each planet, we find $e_{\mathrm{min}} = 0.085_{-0.038}^{+0.047}$ for Kepler-51b and $e_{\mathrm{min}} = 0.059_{-0.037}^{+0.049}$ for Kepler-51d. As illustrated in \autoref{fig:pe_emin_diagnostic}, these photometric requirements are severely incompatible with the robust dynamical eccentricities derived from N-body models \citep[e.g.][]{Masuda2024}, which restrict the eccentricities to $e_b \approx 0.016$ and $e_d \approx 0.01-0.02$. Explaining the anomalies solely through the photoeccentric effect would require eccentricities roughly $3$ to $5$ times larger than the allowed dynamical values. 

\begin{figure}[t]
\centering
\includegraphics[width=\columnwidth]{pe_emin_diagnostic.png}
\caption{The minimum required eccentricity $e_{\mathrm{min}}$ plotted as a function of the density ratio $\Psi = \rho_{\star,\mathrm{obs}}/\rho_{\star,\mathrm{true}}$ (solid black curve). The dashed curves show the actual eccentricity values obtained from the detailed TTV analysis by \citet{Masuda2024}.}
\label{fig:pe_emin_diagnostic}
\end{figure}

Nevertheless, the eccentricities inferred from the TTV analysis are subject to uncertainties; therefore, there is a small probability that some of these eccentricities yield values of $\Psi$ that are consistent with the \photofit posterior distribution. To quantitatively evaluate this probability, we carried out a Monte Carlo calculation of the expected values of the photoeccentric bias, defined as

$$
\Psi^\mathrm{PE} \equiv \frac{(1+e \sin \omega)^3}{\left(1-e^2\right)^{3 / 2}} \, .
$$

In \autoref{fig:pe_emin_diagnostic}, we present the comparison between the distribution of predicted bias provided the uncertainties in $e$ and $e\sin\omega$ from \citet{Masuda2014} and the measured density anomaly for both planets.

\begin{figure}[t]
\centering
\includegraphics[width=0.8\columnwidth]{pe_anomalies.png}
\caption{Top panels: Forward prediction of the photoeccentric bias $\Psi(e,\omega)$ versus the observed density anomaly $\Psi$. Using dynamical posteriors from \citet{Masuda2024}. Bottom panel: The Multibody Asterodensity Profiling (MAP) analysis which compares the expected inter-planet density ratio $\Theta_{bd}=(\Psi_b/\Psi_d)^{2/3}$ and the observed one.}
\label{fig:pe_anomalies}
\end{figure}

As is evident from the predicted PE bias distributions for both planets, the overlap with the observed density anomalies is limited. To formally quantify this, we calculate the probability that the PE effect can fully account for the observed anomaly by determining the fraction of dynamically stable configurations that produce a synthetic density ratio $\Psi^{\mathrm{PE}}$ equal to or greater than the observed ratio $\Psi$. By marginalizing over the dynamical constraints on eccentricity and assuming a uniform distribution for the argument of periastron $\omega$, we find that the probability of the PE effect explaining the anomaly is only $\approx 1\%$ for \exoplanet{Kepler-51}{b} and $\approx 16\%$ for \exoplanet{Kepler-51}{d}. Furthermore, even under the most conservative and optimistic assumption of perfect orbital alignment ($\omega = \pi/2$, yielding the maximum possible PE bias for a given eccentricity), the theoretical upper limits for these probabilities only reach $\approx 3.3\%$ and $\approx 18.7\%$, respectively. These low probabilities strongly disfavor the photoeccentric effect as the primary cause of the anomalously low transit densities.

A final concern remains before we can fully trust the previous result. The analysis relies on the assumption that the true density is known precisely. To conclusively rule out any bias arising from the adopted reference density $\rho_{\star,\mathrm{true}}$, we employ the Multibody Asterodensity Profiling (MAP) framework \citep{kipping2012novel}. MAP exploits the requirement that both planets orbit the same host star. The ratio of their transit-derived densities yields the observable quantity $\Theta_{bd} \equiv (\Psi_b/\Psi_d)^{2/3}$.
Using a first-order expansion of the photoeccentric effect, \citet{kipping2012novel} demonstrated that a deviation of $\Theta_{bd}$ from unity imposes an absolute lower limit on the combined eccentricity of the pair:
\begin{equation}
    (e_b + e_d) \geq \frac{|\Theta_{bd} - 1|}{2} \equiv (e_b + e_d)_{\mathrm{min}}.
\end{equation}

Our analysis yields an observed ratio $\Psi_b/\Psi_d = 1.125_{-0.082}^{+0.090}$, leading to a MAP-inferred lower bound of $(e_b + e_d)_{\mathrm{min}} \approx 0.041$. However, the sum of the dynamical eccentricities from \citet{Masuda2024} is only $(e_b + e_d) \approx 0.024$. The measured eccentricities fall substantially short of the strict multi-planet limit, ruling out the photoeccentric effect independently of the assumed stellar parameters.

\subsection{The \PhotoRing\ Hypothesis}
\label{sec:photoring_hypothesis}

Although the discrepancies could reasonably result from scenarios (ii) or (iii), having rigorously ruled out scenario (i) in \autoref{sec:photoeccentric_effect}, we now turn our attention to scenario (iv). In the remainder of this work we focus on this ringed-planet scenario—not because we consider it the most likely cause of the observed density discrepancies, but in order to evaluate whether a ringed configuration can bring the transit observables into agreement with a self-consistent stellar density. We also emphasize that this is the first study in the literature to perform a full inference using the PR. Therefore, even if other explanations for the density anomalies are possible, using this approach is still a valuable pursuit, as it establishes both a general methodology and computational tools that can be applied to the analysis of other systems. %RF 

If the density anomalies are caused by the presence of rings, a notable feature of the specific case under consideration is that the sign of the measured offset in all our analyses falls in the positive-PR regime ($\rho_{\star,\mathrm{obs}} > \rho_{\star,\mathrm{true}}$), which is less frequently encountered than the negative-$PR$ regime expected for most ring orientations (see \autoref{fig:photoring_contour}). While this does not rule out the ring hypothesis—since physically plausible configurations include both zero and positive offsets—it does indicate that, should rings be driving the observed anomaly, they must occupy a relatively constrained region of parameter space. In particular, this would require the rings to have relatively large tilt angles, $\theta_R \gtrsim 60^\circ$, with respect to the orbital plane. Such large tilts are physically feasible only if the ring inclination relative to the orbital plane is also substantial. The consequences of these configurations will be examined in a later section. %RF

\section{Bayesian Inference: From Observables to Ring Parameters}
\label{sec:bayes_inference}

The forward model introduced in \autoref{sec:geom_model} specifies a deterministic relationship between the ring configuration parameters $\boldsymbol{\theta} = (p, f_i, f_e, i_R, \theta_R, \tau)$—given a set of nuisance parameters $\boldsymbol{\eta} = (P, b, \rho_{\star,\mathrm{true}})$—and the resulting effective transit observables: %RF

\begin{equation}
    \boldsymbol{\theta} \mapsto \mathbf{y}(\boldsymbol{\theta},\boldsymbol{\eta}) \equiv (\predictables).
    \label{eq:effective_observables}
\end{equation}

Our aim is to identify which ring configurations are compatible with the constraints set by the observed transits, while rigorously accounting for both observational and stellar uncertainties in the Kepler-51 system. To achieve this, we must specify the probability distribution function of the predicted parameters $(\predictables)$, conditioned on the available photometric data. %RF

\subsection{Likelihood}
\label{sec:kde_likelihood}

The most natural choice for the pdf of the transit observables is the joint posterior distribution of the photometric observables $(\observables)$. Because the photometric analysis does not yield a continuous distribution, we construct a non-parametric approximation to this function using a multivariate Gaussian kernel density estimator (KDE; \citealt{Scott1992}). To do so, we select a statistically representative subsample of the \photofit-derived samples of the observable vector. We denote this subsample by $\{\mathbf{y}^\mathrm{ph}_m\}_{m=1}^{N}$, where each $\mathbf{y}_m$ corresponds to one joint realization of the chosen observables. These samples therefore provide a discrete characterization of the joint posterior distribution of $\mathbf{y}$ implied by the \photofit, which we approximate with the following multivariate Gaussian KDE: %RF

\begin{equation}
\widehat{p}(\mathbf{y}) = \frac{1}{N}\sum_{m=1}^{N}\,\mathcal{N}\!\left(\mathbf{y}\,\middle|\,\mathbf{y}^\mathrm{ph}_m, \mathbf{H}\right),
\label{eq:kde_def}
\end{equation}
where $\mathbf{H}=h^2\mathbf{\Sigma}$ is the KDE bandwidth matrix, constructed from the sample covariance matrix $\mathbf{\Sigma}$ and a bandwidth scaling factor $h$, which determines the degree of smoothing and the correlations of the kernel in the transit observable space. We adopt the standard automated bandwidth selection implemented in the KDE routine, following Scott's or Silverman's rule \citep{Scott1992, Silverman1986}. 

The resulting continuous function yields an estimate of the probability density across the region of transit observables preferred by the photometric analysis. For any candidate ring configuration $\boldsymbol{\theta}_k$, the forward model generates an associated observable vector, $\mathbf{y}(\boldsymbol{\theta}_k,\boldsymbol{\eta})$ (for the treatment of the nuisance parameters see \autoref{sec:bayes_nuisance} below). We then assess the agreement between this prediction and the photometric posterior by evaluating the KDE density at that point and adopting this value as the likelihood. In this way, correlated uncertainties among the observables are naturally accounted for, without depending on independent point estimates or imposing a specific parametric form for the likelihood:

\begin{equation}
\mathcal{L}(\boldsymbol{\theta}_k) \propto \widehat{p}\!\left(\mathbf{y}=\mathbf{y}(\boldsymbol{\theta}_k,\boldsymbol{\eta}\right).
\label{eq:kde_like}
\end{equation}

The retrieval process consists of sampling the posterior distribution of ring parameters, $\boldsymbol{\theta}$, exploring the region of ring-parameter space whose predicted observables are most consistent with the distribution inferred from the photometric analysis. 

\subsection{Treatment of Nuisance Parameters}
\label{sec:bayes_nuisance}

The orbital period $P$, impact parameter $b$, and true stellar density $\rho_{\star,\mathrm{true}}$ determine the transit chord geometry and orbital scale via the inferred values of $(a/R_\star)_\mathrm{geo}$ and $i_{\rm orb}$ (\autoref{eq:aRstar_true} and \autoref{eq:iorb}). They therefore influence how a given ring configuration translates into the effective transit observables. Because these parameters are not intrinsic to the ring geometry and are instead affected by measurement uncertainties, we treat them as nuisance parameters. Across our retrieval suite, we adopt several different handling strategies for them, ensuring that their uncertainties are propagated into the derived ring properties. %RF

In all setups, $P$ is treated as fixed and set to the literature value (in particular we use the values reported by \citealt{Masuda2024}). Depending on the specific {\bf\em run configuration}, $b$ and/or $\rho_{\star,\mathrm{true}}$ are either held at fiducial values, explicitly sampled, or effectively marginalized over within the likelihood. When these parameters are sampled, we use a truncated normal prior for the impact parameter, $b\sim\mathrm{TruncNorm}(\mu_b,\sigma_b;[0,1])$, centered, again, on the ringless-fit value reported by \citet{Masuda2024}, and a one-dimensional KDE prior for $\rho_{\star,\mathrm{true}}$ based on the isochrone-informed stellar-density posterior from \citet{Berger2023}. %RF

In the most basic setup, when $b$ and/or $\rho_{\star,\mathrm{true}}$ are not treated as sampled parameters, we hold them fixed at their mean values. As a robustness check, we additionally use a pseudo-marginal likelihood in which these parameters are integrated out via Monte Carlo averaging within the likelihood, %RF

\begin{equation}
\mathcal{L}(\boldsymbol{\theta_k}) \approx \frac{1}{M}\sum_{j=1}^{M}\widehat{p}\!\left(\mathbf{y}=\mathbf{y}(\boldsymbol{\theta_k},\boldsymbol{\eta}^{(j)})\right),
\label{eq:pseudo_marg}
\end{equation}
where $\boldsymbol{\theta}$ represents the remaining ring parameters and the random realizations $\boldsymbol{\eta}^{(j)}=(P, \rho_{\star,\mathrm{true}}^{(j)},b^{(j)})$ are sampled from their assumed prior distributions. This method incorporates the external uncertainties in $b$ and $\rho_{\star,\mathrm{true}}$ while retaining the original prior information, without embedding these parameters directly into the sampling chain. This corresponds to the proper Bayesian treatment when such parameters are external constraints that should not be updated by the transit data. %RF

\subsection{Parameter Space and Priors}
\label{sec:bayes_priors}

In all the cases, the retrieval scheme surveys the parameter space $(\parameters)$, keeping the inner ring radius fixed at $f_i=1$, because the current photometric precision is insufficient to resolve structures lying inside the ring’s outer edge. The planetary radius ratio $p$ is either allowed to vary freely or held fixed, in order to accommodate the possibility that a portion of the measured transit depth is contributed by the rings themselves. This would imply a smaller planetary radius and could partially account for the unusually low bulk densities inferred for these super-puff planets when no rings are assumed. In the fixed-radius runs, we set $p=p_{\min}$, corresponding to the limiting–density configuration for the planet (see below). %RF

For the planetary and ring parameters, we adopt priors that lie within physically justified intervals. For the ring's outer radius, we adopt a uniform prior $f_e\sim\mathcal{U}(1,f_{e,\max})$. The upper bound $f_{e,\max}=10$ is chosen based on the stability limit for prograde rings, which can reach roughly $\sim0.5$ times the planetary Hill radius \citep{Domingos2006}. For the ring inclination, we assume an isotropic prior, $p(i_R)\propto\sin i_R$ for $i_R\in[0^\circ,90^\circ]$, and we adopt a uniform prior for the projected tilt, $\theta_R\sim\mathcal{U}(0^\circ,90^\circ)$, which spans all possible ring orientations. %RF

For the planetary radius, we assume a uniform prior on $p$ limited by physically motivated bounds, $p_{\min} \le p \le p_{\mathrm{obs}}$. Here, $p_{\mathrm{obs}}$ is the transit-derived radius under a ringless model (as reported by \citealt{Masuda2024}), and $p_{\min}$ is a lower bound set by requiring that the bulk density does not exceed that of a rocky, Earth-like composition, such that $p_{\min} = (R_\oplus/R_\star)\,(M_p/M_\oplus)^{1/3}$. We adopt the so-called {\em Outside 2:1} solution of \citet{Masuda2024} (their Table~6), taking $M_{p} = 6.9\,M_\oplus$ for both \exoplanet{Kepler-51}{b} and \exoplanet{Kepler-51}{d} (the reported mass ratios $m$). With $R_\star = 0.869\,R_\odot$ \citep{Berger2023}, this yields $p_{\min} = 0.0201$ for both planets (equivalent to $1.90\,R_\oplus$).%RF

When we allow $\tau$ to vary freely, we impose a log-uniform prior $p(\tau)\propto 1/\tau$ over the interval $\tau\in[0.1,10]$. For exploratory analyses, we instead fix $\tau=1$, which represents a ring system with moderate opacity and is similar to the optical depths observed in portions of Saturn’s main rings \citep{Xystouris2023}. This choice prevents the introduction of an extra parameter that is partially degenerate with the ring size and orientation, thereby lowering the dimensionality and computational expense of the retrieval while still providing a representative reference ring model. %RF

For completeness and to facilitate reproducibility, \autoref{tab:priors} reports the priors adopted in our inference procedure. %RF

\begin{table}[t]
\centering
\footnotesize
\caption{Prior distributions used in the ring retrieval. The planetary mass and stellar radius employed to calculate $p_\mathrm{min}$ are taken from \citet{Masuda2024} and \citet{Berger2023}, respectively.}
\label{tab:priors}
\setlength{\tabcolsep}{3pt} % Default is 6pt
\renewcommand{\arraystretch}{1.3}
\begin{tabular}{lll}
\hline
\hline
Parameter & Prior & Range \\
\hline
\multicolumn{3}{l}{\textit{Ring geometry}} \\
$f_e$ & $\mathcal{U}(1, f_{e,\mathrm{max}})$ & $[1, 10]\,R_p$ \\
$f_i$ & Fixed & $f_i = 1$ \\
$i_R$ & $p(i_R)\propto\sin i_R$ & $[0^\circ, 90^\circ]$ \\
$\theta_R$ & $\mathcal{U}(0^\circ, 90^\circ)$ & $[0^\circ, 90^\circ]$ \\
$\tau^{\dagger}$ & $p(\tau)\propto 1/\tau$ & $[0.1, 10]$ \\
\hline
\multicolumn{3}{l}{\textit{Planetary radius}$^\dagger$} \\
$p$ & $\mathcal{U}(p_\mathrm{min}, p_\mathrm{obs})$ & $[0.0201, 0.0722]$ (b) \\
& & $[0.0201, 0.0986]$ (d) \\
\hline
\multicolumn{3}{l}{\textit{Nuisance parameters}$^\dagger$\footnote{Parameters indicated with a $\dagger$ are kept fixed in the baseline configurations and are only varied in the fully unconstrained model.%RF
}} \\
$\rho_{\star,\mathrm{true}}$ & \citet{Berger2023} & $1.896^{+0.218}_{-0.246}$ \\
$b$ & \citet{Masuda2024}  & $0.074\pm{0.072}$ (b)\\
& & $0.003\pm0.095$ (d) \\
\hline
\multicolumn{3}{l}{\textit{Fixed orbital parameters}} \\
$P$ & Fixed \citep{Masuda2024} & $45.154\,\mathrm{b}$ \\
& & $130.186\,\mathrm{d}$ \\
\hline
\hline
\end{tabular}
\end{table}

\subsection{Posterior sampling}
\label{sec:posterior_sampling}

We investigate the posterior distribution of the ring parameters using complementary Bayesian sampling techniques. Our main inference framework relies on nested sampling implemented via the \texttt{dynesty} Python package \citep{Speagle2020}, which yields both posterior samples and estimates of the Bayesian evidence. This method is particularly well suited to handling the complex likelihood structure resulting from the KDE-based observational constraints. %RF

We use a run configuration of $N_{\rm live} = 1200$ for live points, $\Delta\ln Z = 0.01$ as stop condition, and $N_{\rm KDE}=5000$ photometric posterior samples to train the likelihood. To validate the robustness of the inferred posterior structure, we additionally analyze selected configurations using the affine-invariant ensemble MCMC sampler \texttt{emcee} \citep{ForemanMackey2013}. The agreement between both methods provides an independent consistency attention on ring parameter constraints and exploration mode.

%%%%%%%%%%%%%%%%%%%%%%%%%%%%%%%%%%%%%%%%%%%%%%%%%%%%%%%%%%%%%%%%%%%%%%%

\section{Inferred Ring Geometries} 
\label{sec:ring_results}

Given the high dimensionality of the problem (6+3 dimensions, counting the nuisance parameters) and the computational cost of evaluating the KDE over the complete set of observables (another 6 dimensions), the retrieval procedure can quickly become prohibitively slow. Moreover, certain observables are easier to infer from photometric data than others—for instance, determining $\delta$ is significantly simpler than determining $T_{23}$—and parameters like $b$ play a dual role, appearing both as nuisance parameters and as predicted observables. %RF

Consequently, in practice we limited the ring parameters to what we termed the core parameters ($f_e$, $i_R$, and $\theta_R$), which are free in every run, and ($p$, $\tau$), which may be free or fixed and whose treatment defines the {\bf\em input selection}. Rather than collapsing this choice into a single fixed-versus-free dichotomy, we toggle $p$ and $\tau$ {\em independently}: each may be held fixed ($p=p_\mathrm{min}$ and/or $\tau=1$) or sampled, yielding four input selections. In every case we set $f_i=1$, because the available photometric precision is insufficient to resolve structures inside the outer edge of the ring. In parallel, distinct {\bf\em observable sets} were analyzed, each defined by a different subset of observables: $\mathbf{y}=(\delta,\rho_{\star,\mathrm{obs}})$, $\mathbf{y}=(\delta,\rho_{\star,\mathrm{obs}},T_{14})$, and $\mathbf{y}=(\delta,T_{14},T_{23})$. %RF

\subsection{Results of retrieval setups}
\label{sec:retrieval_setups}

In summary, we categorize our retrieval experiments along three axes: (1)~{\bf\em run configuration}---the $2\times2$ grid of whether the nuisance parameters $\rho_{\star,\mathrm{true}}$ and $b$ are held fixed or sampled (four options: both fixed; only $\rho_{\star,\mathrm{true}}$ free; only $b$ free; both free); (2)~{\bf\em input selection}---the analogous $2\times2$ grid for the free inputs $p$ and $\tau$ (four options); and (3)~{\bf\em observable set}---which subset of transit observables enters the KDE likelihood (three options, listed above). A {\em\bf retrieval setup} is one point in this product space. Counting both planets, the full suite therefore comprises
$2\,\mathrm{(planets)}\times 3\,\mathrm{(observable~sets)}\times 4\,\mathrm{(run~configurations)}\times 4\,\mathrm{(input~selections)}=96$
independent nested-sampling runs. %RF

Restricted retrieval setups allow us to disentangle clear degeneracies among ring parameters such as ring size, orientation, and opacity. For instance, distinct ring configurations can produce nearly identical transit depths while leading to different transit durations or stellar-density anomalies. Indeed, even a stellar density measurement that appears unbiased can still be consistent with particular ring setups (see the zero level in \autoref{fig:photoring_contour}). Thus, combining multiple observables is crucial for breaking these degeneracies and obtaining more robust constraints on the ring properties. Another interesting situation where it is helpful to experiment with different retrieval setups concerns the planetary radius. For example, by fixing this parameter to $p_\mathrm{min}$, we can explore which ring configurations both match the observed transit depth—without requiring an inflated planet—and reproduce the measured density anomaly. %RF

In \autoref{tab:ring_summary}, we present the results of the best 12 retrieval setups for both \exoplanet{Kepler-51}{b} and \exoplanet{Kepler-51}{d}. The tabulated solutions are drawn from the full 96-run suite described above. Because the robustness of the inferred posteriors is highly sensitive to the adopted retrieval setup, we select preferred solutions according to two main criteria: (i) the convergence behavior and shape of the marginal posteriors, favoring unimodal, tightly constrained distributions, and (ii) the capacity of the ring models to match the photometric and photodynamical observables in posterior predictive checks, including those observables that were omitted from the likelihood. %RF

As an additional derived parameter, we report in the final column the bulk planetary density $\rho_p$ of the ringed planet, inferred from $p$ under the assumption of planetary masses
$M_b = M_d = 6.9\,M_\oplus$ \citep{Masuda2024}, and adopting
$R_\star = 0.869\,R_\odot$ \citep{Berger2023}. For comparison, the planetary densities obtained directly from the photodynamical modeling under the assumption of no rings are $0.11$ g cm$^{-3}$ for \exoplanet{Kepler-51}{b} and $0.038$ g cm$^{-3}$ for \exoplanet{Kepler-51}{d}. %RF

As expected, in all ringed-retrieval setups, the inferred planetary density exceeds that of the no-rings photometric fit. Nevertheless, with the exception of the cases where a minimum planetary radius is imposed, whenever the planetary radius is allowed to vary freely in the inference procedure, the resulting density still falls short of the typical values for planets in this mass range. Thus, according to our PR-based analysis of the stellar-density anomaly, both objects remain low-density planets. This indicates that, at least in the case of Kepler-51, the ring hypothesis does not fully resolve the low-density puzzle, but instead points toward a distinct planetary class: puff ringed planets.

\begin{table*}[t!]
\centering
\caption{Optimal ring parameters for selected PR-retrieval configurations for \exoplanet{Kepler-51}{b} and d. Reported values are the posterior medians with 68\% credible intervals. Setups are ordered by planet, by the chosen observable set (i.e. which observables are targeted and define the KDE), and by run configuration (i.e. which nuisance or input parameters are fixed versus free).}
\label{tab:ring_summary}
\scriptsize
\setlength{\tabcolsep}{0pt}
%%% GENERATED BY CODE
% run_version: 2026080215
\begin{tabular*}{\textwidth}{@{\extracolsep{\fill}} c @{\hspace{0.3em}} c @{\hspace{0.3em}} c @{\hspace{0.3em}} c @{\hspace{0.3em}} c @{\hspace{0.3em}} c @{\hspace{0.3em}} c @{\hspace{0.3em}} c @{\hspace{0.3em}} c @{\hspace{0.3em}} c @{\hspace{0.3em}} c @{}}
\toprule
& Observable set & Run config. & \multicolumn{7}{c}{Input selection} & Other\\
\cmidrule(lr){4-10}
Planet &  ($\mathcal{L}_{\rm KDE}$) & Nuis.+Input
  & $f_e\;[R_p]$ & $i_R\;[^\circ]$ & $\theta_R\;[^\circ]$
  & $p\;[R_\star]\;(R_\oplus)$ & $\tau$
  & $\rho_{\star,\rm true}\;[\rm g\,cm^{-3}]$ & $b$ & $\rho_p\;[\rm g\,cm^{-3}]$ \\
\midrule

% ======================== PLANET B ========================
\multirow{7}{*}{\rotatebox{90}{\exoplanet{Kepler-51}{b}}}
  & \multirow{3}{*}{$\delta, \rho_{\star,\rm obs}$}
    & all-free
    & $2.70_{-1.09}^{+2.26}$
    & $73.82_{-17.55}^{+9.54}$
    & $77.68_{-10.45}^{+8.97}$
    & $0.05_{-0.02}^{+0.01}\;(4.97_{-2.12}^{+1.18})$
    & $1.22_{-0.94}^{+2.77}$
    & $1.96_{-0.23}^{+0.26}$
    & $0.14_{-0.10}^{+0.11}$
    & $0.31_{-0.15}^{+1.33}$ \\

  & 
    & $\rho_{\star,\rm true}$
    & $5.43_{-0.47}^{+1.36}$
    & $61.95_{-10.01}^{+12.11}$
    & $73.20_{-5.95}^{+11.17}$
    & $0.0201^\dagger\;(1.90^\dagger)$
    & $1.00^\dagger$
    & $1.97_{-0.22}^{+0.19}$
    & $0.0740^\dagger$
    & $5.51$ \\

  & 
    & all-fixed
    & $5.38_{-0.30}^{+1.04}$
    & $61.06_{-5.83}^{+10.78}$
    & $73.47_{-3.90}^{+10.25}$
    & $0.0201^\dagger\;(1.90^\dagger)$
    & $1.00^\dagger$
    & $1.883^\dagger$
    & $0.0740^\dagger$
    & $5.51$ \\
\cmidrule{2-11}


  & \multirow{3}{*}{$\delta, T_{14}, \rho_{\star,\rm obs}$}
    & all-free
    & $1.63_{-0.51}^{+4.94}$
    & $65.77_{-24.14}^{+13.71}$
    & $71.65_{-14.59}^{+10.91}$
    & $0.064_{-0.039}^{+0.006}\;(6.09_{-3.70}^{+0.54})$
    & $0.55_{-0.36}^{+1.31}$
    & $2.31_{-0.19}^{+0.11}$
    & $0.23_{-0.08}^{+0.09}$
    & $0.17_{-0.04}^{+2.62}$ \\

  & 
    & $\rho_{\star,\rm true}+b+p$
    & $1.27_{-0.22}^{+1.33}$
    & $49.18_{-19.88}^{+16.66}$
    & $69.28_{-18.38}^{+14.15}$
    & $0.066_{-0.023}^{+0.004}\;(6.27_{-2.19}^{+0.36})$
    & $1.00^\dagger$
    & $2.37_{-0.18}^{+0.12}$
    & $0.19_{-0.13}^{+0.11}$
    & $0.15_{-0.02}^{+0.40}$ \\

  & 
    & $\rho_{\star,\rm true}+b$
    & $5.20_{-0.43}^{+1.80}$
    & $57.89_{-11.99}^{+17.19}$
    & $73.14_{-8.22}^{+11.46}$
    & $0.0201^\dagger\;(1.90^\dagger)$
    & $1.00^\dagger$
    & $2.26_{-0.17}^{+0.16}$
    & $0.26_{-0.10}^{+0.08}$
    & $5.51$ \\
\cmidrule{2-11}


  & $\delta, T_{14}, T_{23}$
    & all-free
    & $1.51_{-0.38}^{+6.59}$
    & $73.91_{-23.77}^{+11.40}$
    & $76.44_{-14.68}^{+8.76}$
    & $0.067_{-0.040}^{+0.003}\;(6.37_{-3.82}^{+0.27})$
    & $0.35_{-0.21}^{+1.71}$
    & $2.35_{-0.14}^{+0.09}$
    & $0.21_{-0.07}^{+0.08}$
    & $0.15_{-0.02}^{+2.13}$ \\
\midrule

% ======================== PLANET D ========================
\multirow{6}{*}{\rotatebox{90}{\exoplanet{Kepler-51}{d}}}
  & \multirow{2}{*}{$\delta, \rho_{\star,\rm obs}$}
    & all-free
    & $1.80_{-0.62}^{+1.68}$
    & $64.25_{-27.13}^{+18.23}$
    & $68.30_{-17.70}^{+14.39}$
    & $0.083_{-0.045}^{+0.011}\;(7.90_{-4.27}^{+1.07})$
    & $1.76_{-1.45}^{+4.27}$
    & $2.03_{-0.19}^{+0.15}$
    & $0.14_{-0.09}^{+0.13}$
    & $0.077_{-0.024}^{+0.724}$ \\

  & 
    & all-fixed
    & $6.95_{-0.31}^{+1.01}$
    & $56.17_{-6.91}^{+11.27}$
    & $70.25_{-5.29}^{+12.08}$
    & $0.0201^\dagger\;(1.90^\dagger)$
    & $1.00^\dagger$
    & $1.883^\dagger$
    & $0.0030^\dagger$
    & $5.51$ \\
\cmidrule{2-11}


  & \multirow{3}{*}{$\delta, T_{14}, \rho_{\star,\rm obs}$}
    & all-free
    & $1.64_{-0.47}^{+3.86}$
    & $52.53_{-24.65}^{+21.84}$
    & $64.36_{-12.85}^{+13.80}$
    & $0.075_{-0.045}^{+0.017}\;(7.10_{-4.23}^{+1.63})$
    & $1.51_{-0.99}^{+2.09}$
    & $2.16_{-0.14}^{+0.18}$
    & $0.27_{-0.14}^{+0.07}$
    & $0.11_{-0.05}^{+1.50}$ \\

  & 
    & $\rho_{\star,\rm true}+b+p$
    & $1.26_{-0.20}^{+0.63}$
    & $63.89_{-21.37}^{+18.86}$
    & $59.08_{-33.42}^{+18.25}$
    & $0.09_{-0.01}^{+0.00}\;(8.76_{-1.33}^{+0.25})$
    & $1.00^\dagger$
    & $2.17_{-0.12}^{+0.15}$
    & $0.27_{-0.09}^{+0.06}$
    & $0.057_{-0.005}^{+0.036}$ \\

  & 
    & $\rho_{\star,\rm true}+p$
    & $2.15_{-0.68}^{+4.38}$
    & $41.62_{-14.87}^{+21.85}$
    & $55.73_{-16.41}^{+9.18}$
    & $0.06_{-0.03}^{+0.02}\;(5.43_{-3.24}^{+1.81})$
    & $1.00^\dagger$
    & $2.427_{-0.034}^{+0.032}$
    & $0.0030^\dagger$
    & $0.24_{-0.14}^{+3.38}$ \\
\cmidrule{2-11}


  & $\delta, T_{14}, T_{23}$
    & all-free
    & $2.77_{-1.37}^{+3.25}$
    & $68.24_{-14.97}^{+8.89}$
    & $69.02_{-12.56}^{+14.74}$
    & $0.06_{-0.03}^{+0.04}\;(5.33_{-2.72}^{+3.41})$
    & $0.85_{-0.56}^{+2.26}$
    & $2.16_{-0.21}^{+0.21}$
    & $0.28_{-0.15}^{+0.09}$
    & $0.25_{-0.19}^{+1.88}$ \\
\bottomrule
\end{tabular*}

%%% GENERATED BY CODE
\vspace{1ex}
{\raggedright \textit{Note:} The parameters that may vary within each run configuration (free) include the nuisance parameters, $\rho_\mathrm{\star,true}$ and $b$, as well as the input parameters, $p$ and $\tau$. Whenever one of these parameters is held fixed, its value is indicated with a $\dagger$.\par}
\end{table*}

\subsection{Marginal Posteriors}
\label{sec:ring_marginals}

\paragraph{Ring outer radius $f_e$}

Across all preferred configurations, the obtained posterior of $f_e$ is well-defined and away from both uniform prior boundaries. For \exoplanet{Kepler-51}{b}, the median converges to $f_e\approx1.95$--$2.71\,R_p$ depending on the likelihood choice, while for \exoplanet{Kepler-51}{d} it settles at $f_e\approx1.69$--$2.07\,R_p$. These values correspond to compact ring systems, well below the stability limit of $\sim0.5,R_{\rm Hill}$ ($f_{e,\max}\sim10\,R_p$; \citealt{Domingos2006}) where stellar tides remain subdominant. At the orbital distances of \exoplanet{Kepler-51}{b} ($P \approx 45\,$d) and d ($P \approx 130\,$d) tidal forces are too weak to cause significant planetary deformation, disruption or hydrodynamic atmospheric escape—effects that only occur at much shorter orbital distances ($P<10\,$d). This implies that our inferred ring solutions survive unaffected by stellar tidal effects. For comparison, Saturn's main rings extend to $\approx 2.3\,R_{p}$, demonstrating that the inferred ring sizes are physically plausible, at least within a Solar System context.

The most persistent pattern across the full retrieval suite is the systematic anti-correlation between $f_e$ and the planetary radius ratio $p$. When $p$ is held fixed at $p_{\min}$, the posterior medians of $f_e$ consistently converge to the range $4$--$7\,R_p$ for both planets across all likelihood configurations; when $p$ is set free, $f_e$ falls to $\lesssim3\,R_p$. This behaviour is a consequence of the structure of the forward model: the observed transit depth $\delta$ is proportional to the effective occulting area $A_{\rm eff}=\pi R_p^2 + A_{\rm ring}$, so any combination of planetary radius and ring area that preserves $A_{\rm eff}$ produces an identical $\delta$. A likelihood built on $\delta$ alone cannot distinguish between these combinations, and the sampler freely explores the iso-depth surface in $(f_e, p)$ space.

%  The observational implication is that Kepler long-cadence transit photometry, by itself, cannot simultaneously determine the true planetary size and the ring extent: an independent, externally motivated prior on $p$—for instance, one derived from TTV orbital dynamics or from multi-wavelength transit observations sensitive to the atmospheric opacity—is required to break this fundamental degeneracy.}

\paragraph{Ring orientation: $i_R$ and $\theta_R$}

The posterior of the projected ring inclination, $i_R$, remains relatively broad, with median values consistently exceeding $50^\circ$ for both planets. Because the adopted isotropic prior, $p(i_R)\propto\sin i_R$, already favors edge-on configurations, any departure from the prior provides a measure of the information contributed by the data. In the preferred retrievals, the posteriors exhibit a clear preference for moderately inclined rings ($i_R\approx60^\circ$--$75^\circ$), whereas in less informative configurations (fewest free parameters) they largely retain the prior shape, indicating weaker constraints on the ring inclination. Such substantial obliquity is physically expected for the young Kepler-51 system ($\sim500$ Myr). While planets typically form with low obliquities under standard formation scenarios, the timescales available for secular resonances and giant impacts are sufficient to tilt the planetary spin axes. This evolutionary history is supported by the large TTVs observed in the system, which testify to ongoing dynamical interactions among the planets.

In contrast, the projected azimuthal tilt, $\theta_R$, is more tightly constrained, with posterior medians concentrated around $\theta_R\approx60^\circ$--$75^\circ$ for planet b and $\theta_R\approx60^\circ$--$68^\circ$ for planet d. This behavior reflects an intrinsic $i_R$--$\theta_R$ correlation present in all likelihood configurations, since the observed transit depth $\delta$ always requires a balance between ring inclination and projected area, and including $T_{14}$ in the likelihood restricts this degeneracy because the same orientation must simultaneously reproduce the contact geometry with the stellar limb. A key internal consistency attention between the analytical PR framework and the Bayesian retrievals arise in configurations where $\rho_{\star,\mathrm{obs}}$ and $b$ are also free sampled; these preferred ring orientations lie predominantly within the positive-$PR$ regime ($\rho_{\star,\rm obs}>\rho_{\star,\rm true}$; \autoref{fig:photoring_contour}), in agreement with the sign of the density anomaly independently inferred from our \photofit for both Kepler-51 planets (Fig.~\ref{fig:rhoobs_comparison}). The recovery of orientations within this restricted region, without imposing any prior preference on the sign of the anomaly, demonstrates that the likelihood constructed from Kepler data naturally favors the ring geometries required to reproduce the observed stellar-density discrepancy.

\paragraph{Planetary radius $p$}

When treated as a free parameter, $p$ converges to a narrow range: $p\approx0.045$--$0.065$ for planet b and $p\approx0.07$--$0.09$ for planet d. These values are smaller than the ringless-model estimates of $p_{\rm obs}=0.072$ (b) and $p_{\rm obs}\approx0.098$ (d) from \citet{Masuda2024}. The resulting bulk densities are therefore enhanced by factors of a few, reaching at most $\sim4$ for planet b and $\sim3$ for planet d, partially reducing, without entirely removing, their extreme super-puff nature. While rings are not required to be the sole explanation for the anomalously low densities of the Kepler-51 planets, our retrievals demonstrate that a self-consistent ring geometry exists in which a significant fraction of the apparent radius inflation can be interpreted as a photometric effect arising from circumplanetary material.

In the fixed-$p_{\min}$ setup—the lower physical bound corresponding to a maximum-density rocky body—the solutions require extended rings whose projected area dominates the transit signal, implying ring-to-planet area ratios far exceeding those observed in the Solar System. These configurations tend to overpredict the transit durations $T_{14}$ and $T_{23}$, revealing a geometric tension that can only be solved by introducing additional freedom in $b$ and/or $\rho_{\star,\rm true}$, or by explicitly incorporating the transit durations into the likelihood to restrict the extended ring solutions. The fixed-$p_{\min}$ retrievals are therefore best interpreted as a limiting case for maximum ring extent compatible with the data, rather than a physically favored explanation.

\subsection{Effects of Nuisance Parameters}
\label{sec:ring_nuisance}

The most physically important result of the nuisance parameter treatment is that the preferred ring geometry---and the ability to simultaneously reproduce all five photometric-derived transit observables---is recovered only when both $\rho_{\star,\mathrm{true}}$ and $b$ are sampled jointly as free parameters. Because these two quantities defines the orbital scale and the transit chord, respectively, their coupled freedom allows the forward model to balance the ring contact geometry against the transit durations without forcing ring parameters to absorb a geometric tension it cannot physically explain.  We discuss the behavior of each parameter in turn before summarizing their joint effect.

\paragraph{True stellar density $\rho_{\star,\rm true}$.}

In the \emph{fixed} $\rho_{\star,\rm true}$ setup, adopting the mean value of $1.8832\,\rm g\,cm^{-3}$ from \citet{Berger2023} distribution, the forward model forces the ring geometry to absorb the full tension between the fixed orbital scale and the \photofit-derived observables. We obtain informative ring posteriors, far from the priors, only when the likelihood is restricted to $(\delta, \rho_{\star,\rm obs})$. In all other likelihood configurations, the recovered posteriors largely retain the shape of the priors, indicating minimal constraining power of data under a fixed orbital scale. Furthermore, although the model reproduces $\delta$ and $\rho_{\star,\rm obs}$ when these quantities are explicitly included in the likelihood, it fails to predict the remaining observables, particularly transit durations, revealing that the fixed orbital scale is inconsistent with the ring-extended contact geometry required to match the duration.

Conversely, while the pseudo-marginalisation strategy yields well-converged posteriors for the ring parameters, consistent with the expected positive-PR orientation region, but the resulting configurations fail to reproduce the entire transit observables set. By averaging the likelihood over the full $\rho_{\star,\rm true}$ distribution at each evaluation, this approach prioritizes statistical compatibility with stellar density uncertainty over fitting the transit data. Consequently, pseudo-marginalisation improves methodological consistency when characterizing plausible ring configurations, but it sacrifices the predictive precision required to identify the specific geometry that best reproduces the Kepler light curves.

When sampled as a free parameter, the posterior, depends critically on the likelihood configuration. Whit only $(\delta, \rho_{\star,\rm obs})$ in the likelihood, the posterior preserves the \citet{Berger2023} prior shape, converging to medians of $\rho_{\star,\rm true} \approx 1.88$--$1.99\,\rm g\,cm^{-3}$ for both planets, fully consistent with those independent isochrone-based estimates in \autoref{tab:rho_star_compendium}. In this setup the transit depth and observed stellar density can be jointly reproduced without modifying of the orbital scale, but the transit geometry remains effectively unconstrained, so
the transit durations and impact parameter are poorly reproduced. Once the full transit duration, $T_{14}$, is added to the likelihood, the $\rho_{\star,\mathrm{true}}$ posterior shifts toward higher values ($\rho_{\star,\rm true} \approx 2.05$--$2.47\,\rm g\,cm^{-3}$) because the sampler exploits a degeneracy in which a higher stellar density compresses the orbital timescale while the ring geometry simultaneously extends the transit duration through the outer contact positions. Although this configuration yields the best posterior predictive performance, the resulting increase in $\rho_{\star,\rm true}$ is physically undesirable, as it becomes inconsistent with the isochrone-based estimations. This is particularly relevant because our framework explicitly distinguishes between the photometric observed density, $\rho_{\star,\rm obs}$, and the independently meausered stellar density, $\rho_{\star,\rm true}$. While $\rho_{\star,\rm obs}$ is expected to be biased under our ring hypothesis, $\rho_{\star,\rm true}$ is an intrinsic stellar property and therefore should not be modified by the ring itself. The upward shift in $\rho_{\star,\rm true}$ should therefore be interpreted as a consequence of the model balancing the orbital scale and ring geometry to reproduce the transit observables, rather than as evidence for a revision of the stellar properties or a new inference on the stellar interior.

\paragraph{Impact parameter $b$}

Fixing $b$ to the JWST-derived reference value \citep[$b=0.0740$ for planet b, $0.0030$ for planet d;][]{Masuda2024} removes the uncertainty in the transit chord and leaves the ring parameters as the only handle on the contact time excess. For \exoplanet{Kepler-51}{d}, whose transit is nearly central,  reproducing a duration excess under this constraint requires an almost edge-on ring orientation and maximal radial extent, driving $f_e\to f_{e,\rm max}$ and $i_R\to90^\circ$ in configurations where $T_{14}$ contributes to the likelihood. However, the likelihood constraint on $\delta$ prevents these limits from being reached, since reproducing the observed depth requires a balance between $f_e$, $p$, and $i_R$, but the model still fails to simultaneously reproduce the contact times without the compensating effect of sampling $\rho_{\star,\rm true}$. Moreover, fixing $b$ neglects the uncertainty in $b_{\rm obs}$ imposed by Kepler long-cadence photometry, resolving a degeneracy that the available data cannot constrain. In contrast, accounting for the uncertainty in $b$ through the pseudo-marginalization over the \citet{Masuda2024} prior reduces the need for extreme edge-on and extended configurations to compensate the excess on transit duration. However, as discussed above in the same treatment for $\rho_{\star,\rm true}$, this methodological improve comes at the cost of a poorer reproduction of the whole transit observables set.

Allowing $b$ to vary, jointly with $\rho_{\star,\mathrm{true}}$, is the condition under which the preferred ring posteriors discussed in the previous section are recovered and the five transit observables are simultaneously reproduced most successfully. When the likelihood is restricted to $(\delta, \rho_{\star,\mathrm{obs}})$, the posterior largely preserves the prior shape reported by \citet{Masuda2024}, reflecting the insufficient contact geometry information, as $\rho_{\star,\mathrm{obs}}$ accounts for the global distortion effect of rings, while $\delta$ primarily constrains the ring's effective area. Moreover, in configurations constrained by a likelihood including $T_{14}$, its posterior converges to median values of $b\approx0.13$--$0.27$ for \exoplanet{Kepler-51}{b} and $b\approx0.14$--$0.34$ for Kepler-51,d, significantly higher than the near-central values inferred by \citet{Masuda2024}. In this configurations, the forward model shifts the transit contacts to higher stellar latitudes with extended ingress and egress compatible with the broad \photofit-derived transit durations. However, the $\sim30$-min cadence of Kepler photometry does not temporally resolve this contact morphology sufficiently to disentangle duration increases produced by rings from those associated with larger impact parameters. In any case, the posterior of $b$ should be interpreted as a model-conditional geometric artifact and does not impy evidence for a revised orbital architecture. Distinguishing between these two effects will require the simultaneous fitting of JWST photometry, which provides the cadence needed to directly resolve ingress/egress slope from ring contact structure.

\subsection{Joint posterior distribution}
\label{sec:ring_predictive}

As should now be evident from the previous sections, identifying the optimal retrieval setup is a nontrivial task. After reviewing all filtered setups summarized in \autoref{tab:ring_summary}, we chose the one that most closely matches the typical best-fit values of the ring parameters. The selected configuration is the setup in which all nuisance and input parameters are allowed to vary freely (``all-free''), and the set of observables is restricted to $\rho_\mathrm{\star,obs}$, $\delta$, and $T_{14}$. 

In \autoref{fig:ring_panel_b} and \autoref{fig:ring_panel_d}, we show corner plots of the joint posterior distributions for both the ringed-planet parameters and the nuisance parameters of \exoplanet{Kepler-51}{b} and \exoplanet{Kepler-51}{d}, corresponding to the chosen retrieval setup. For illustration, we also display a graphical depiction of the planet–ring system, including the sizes and orientations, assuming the median values of the fitted parameters. As can be seen, the preferred sizes and orientations of the planets are consistent with the conclusions reached in the previous sections. Our earlier result that the ring solutions for both planets still imply low-density planets remains valid. Likewise, the relatively large inclination of the ring plane with respect to the planetary equator is also preserved. %RF

\begin{figure*}[t]
\centering
\includegraphics[width=\textwidth,height=\textheight,keepaspectratio]{final_planetb_panel_reduced.png}
\caption{Posterior joint distribution for the PR inference of \exoplanet{Kepler-51}{b} for a chosen retrieval setup (see text). The corner-plot panels display 2D and 1D marginal posteriors for both the ringed planet parameters and the nuisance parameters. For visualization, we include a schematic of the ring system constructed from the median values of the marginal posteriors, indicated by red lines on the 1D marginals. The dark disk denotes the planet, while the filled elliptical annulus represents the ring. Red dots on the panels indicate where the ringed planet parameter is positioned that matches the schematic illustration.%RF
}
\label{fig:ring_panel_b}
\end{figure*}


\begin{figure*}[t]
\centering
% --- \exoplanet{Kepler-51}{d} ---
\includegraphics[width=\textwidth,height=\textheight,keepaspectratio]{final_planetd_panel_reduced.png}

\caption{Posterior joint distribution for the PR inference of \exoplanet{Kepler-51}{d} for a chosen retrieval setup (see text). The figure specifications are identical to those in \autoref{fig:ring_panel_b}. %RF
}
\label{fig:ring_panel_d}
\end{figure*}

\subsection{Predictive Posterior Ckecks}
\label{sec:ring_predictive}

The final test we perform on the retrieval results is to verify whether the best-fit input and nuisance parameters (when allowed to vary) can reproduce the marginal posterior distributions of most observables derived from the photometric analysis, including those that were excluded from the likelihood evaluation. We carry out these goodness-of-fit assessments using posterior predictive checks (PPCs). Specifically, for each retrieval setup, we propagate the corresponding posterior samples through the forward model described in \autoref{sec:geom_model}. We then compare the resulting marginal posterior distributions of the transit observables with those inferred directly from the photometric data (\autoref{sec:data_analysis}). For validation, we also include observables that were not explicitly incorporated in the likelihood calculation. The agreement for each observable is quantified using two metrics (see eg. \citealt{Ramdas2017}): (i) the 1-Wasserstein distance $W_1$, which characterizes the mean absolute shift between the two empirical distributions; and (ii) the energy distance $E$, which quantifies the squared difference between the two distributions. %RF

For illustration in \autoref{fig:ppcs_planet}, we display and compare the marginal posterior distributions (shown as histograms) of the observables $(\delta, T_{14}, T_{23}, \rho_{\star,\mathrm{obs}}, b_{\mathrm{obs}})$ obtained from the photometric analysis with those derived from the retrieval setup whose selected results are presented in \autoref{fig:ring_panel_b} and \autoref{fig:ring_panel_d}. In this example, the ring model closely matches all the observed characteristics of the light curves, which is further confirmed by the corresponding goodness-of-fit metric values. %RF

\begin{figure*}[t]
    \centering
    % \hfill
    % \begin{minipage}[t]{0.42\textwidth}
    %     \centering
    %     \includegraphics[width=0.95\linewidth]{figures/ppc_planet_b_vertical.png}
    % \end{minipage}
    % \hfill
    % \begin{minipage}[t]{0.42\textwidth}
    %     \centering
    %     \includegraphics[width=0.95\linewidth]{figures/ppc_planet_d_vertical.png}
    % \end{minipage}
    % \hfill
    \includegraphics[width=0.5\linewidth]{figures/ppc_planet_b_vertical.png}
    \includegraphics[width=0.5\linewidth]{figures/ppc_planet_d_vertical.png}
    \caption{Posterior predictive checks (PPCs) for \exoplanet{Kepler-51}{b} and \exoplanet{Kepler-51}{d} in the \textit{fully free} retrieval setup. PPCs compare the model-predicted distribution (blue) with the photometrically derived posterior (grey) for each transit observable. The reported values of $W_1$ and $E$ quantify the self-consistent goodness-of-fit between the ringed-planet model and the observed transit observables.%RF
    }
    \label{fig:ppcs_planet}
\end{figure*}

\autoref{tab:ppc_all} presents the values of the PPCs metrics for the selected observable predictions across all the 12 selected retrieval setups.  

\begin{table*}[t]
\centering
\caption{Posterior predictive attention (PPC) statistics for the chosen retrieval setups of \exoplanet{Kepler-51}{b} and d, grouped by planet, observable set and run configuration (same as \autoref{tab:observables_full}). For each observable column, the two values in each cell give the 1-Wasserstein distance $W_1$ and the energy distance $E$ (see main text). \textbf{Bold} entries denote observables that enter $\mathcal{L}_{\rm KDE}$ (in-sample), whereas entries in regular font correspond to out-of-sample predictions used as a predictive validation of the ring model. The impact parameter $b_{\rm obs}$ is always treated as out-of-sample.%RF
}
\label{tab:ppc_all}
\scriptsize
\setlength{\tabcolsep}{0pt}
%%%% GENERATED BY CODE %%%%%
% run_version: 2026080215
\begin{tabular*}{\textwidth}{@{\extracolsep{\fill}} c c c c c c c c @{}}
\toprule
Planet &
Observable set &
Run config. &
$\delta$ [ppm] &
$T_{14}$ [h] &
$T_{23}$ [h] &
$\rho_{\star,\rm obs}$ [g\,cm$^{-3}$] &
$b_{\rm obs}$\\
&  ($\mathcal{L}_{\rm KDE}$) & Nuis.+Input & & & & & \\
\midrule

% ======================== PLANET B ========================
\multirow{7}{*}{\rotatebox{90}{\exoplanet{Kepler-51}{b}}}
  & \multirow{3}{*}{$\delta, \rho_{\star,\rm obs}$}
    & all-free
    & \textbf{5.96\,/\,0.708}
    & 0.372\,/\,0.658
    & 0.417\,/\,0.697
    & \textbf{0.0309\,/\,0.0726}
    & 0.0601\,/\,0.154 \\

  & 
    & $\rho_{\star,\rm true}$
    & \textbf{2.64\,/\,0.297}
    & 0.441\,/\,0.789
    & 0.5\,/\,0.843
    & \textbf{0.0154\,/\,0.0379}
    & 0.0297\,/\,0.0853 \\

  & 
    & all-fixed
    & \textbf{5.32\,/\,0.59}
    & 0.494\,/\,0.974
    & 0.56\,/\,1.04
    & \textbf{0.0141\,/\,0.0332}
    & nan\,/\,nan \\
\cmidrule{2-8}


  & \multirow{3}{*}{$\delta, T_{14}, \rho_{\star,\rm obs}$}
    & all-free
    & \textbf{5.15\,/\,0.542}
    & \textbf{0.0121\,/\,0.0575}
    & 0.00552\,/\,0.0288
    & \textbf{0.0416\,/\,0.106}
    & 0.0367\,/\,0.107 \\

  & 
    & $\rho_{\star,\rm true}+b+p$
    & \textbf{11.1\,/\,1.34}
    & \textbf{0.0101\,/\,0.0486}
    & 0.00843\,/\,0.0425
    & \textbf{0.0195\,/\,0.0517}
    & 0.0185\,/\,0.0533 \\

  & 
    & $\rho_{\star,\rm true}+b$
    & \textbf{3.07\,/\,0.308}
    & \textbf{0.00977\,/\,0.0468}
    & 0.0036\,/\,0.0188
    & \textbf{0.0352\,/\,0.0904}
    & 0.0332\,/\,0.0973 \\
\cmidrule{2-8}


  & $\delta, T_{14}, T_{23}$
    & all-free
    & \textbf{5.58\,/\,0.634}
    & \textbf{0.00657\,/\,0.0323}
    & \textbf{0.00445\,/\,0.0233}
    & 0.0185\,/\,0.0534
    & 0.0222\,/\,0.0686 \\
\midrule

% ======================== PLANET D ========================
\multirow{6}{*}{\rotatebox{90}{\exoplanet{Kepler-51}{d}}}
  & \multirow{2}{*}{$\delta, \rho_{\star,\rm obs}$}
    & all-free
    & \textbf{35\,/\,2.38}
    & 0.383\,/\,0.627
    & 0.435\,/\,0.651
    & \textbf{0.0423\,/\,0.0913}
    & 0.0479\,/\,0.122 \\

  & 
    & all-fixed
    & \textbf{31.3\,/\,2.1}
    & 0.631\,/\,1.08
    & 0.742\,/\,1.17
    & \textbf{0.0374\,/\,0.0759}
    & 0.135\,/\,0.364 \\
\cmidrule{2-8}


  & \multirow{3}{*}{$\delta, T_{14}, \rho_{\star,\rm obs}$}
    & all-free
    & \textbf{3.93\,/\,0.226}
    & \textbf{0.00638\,/\,0.0243}
    & 0.00262\,/\,0.0072
    & \textbf{0.014\,/\,0.0325}
    & 0.0145\,/\,0.0375 \\

  & 
    & $\rho_{\star,\rm true}+b+p$
    & \textbf{9.57\,/\,0.611}
    & \textbf{0.00647\,/\,0.0262}
    & 0.00449\,/\,0.0141
    & \textbf{0.0212\,/\,0.0531}
    & 0.0194\,/\,0.0529 \\

  & 
    & $\rho_{\star,\rm true}+p$
    & \textbf{18.5\,/\,1.23}
    & \textbf{0.0123\,/\,0.0473}
    & 0.00968\,/\,0.0298
    & \textbf{0.0162\,/\,0.0354}
    & 0.0113\,/\,0.0272 \\
\cmidrule{2-8}


  & $\delta, T_{14}, T_{23}$
    & all-free
    & \textbf{20.4\,/\,1.22}
    & \textbf{0.00227\,/\,0.0095}
    & \textbf{0.00947\,/\,0.0255}
    & 0.00902\,/\,0.0187
    & 0.00544\,/\,0.0149 \\
\bottomrule
\end{tabular*}

%%%% GENERATED BY CODE %%%%%
\end{table*}

A key result is that incorporating a particular subset of observables into the likelihood does not ensure that the entire set of transit observables is accurately reproduced. In the most basic setup ($\mathbf{y}=(\delta,\rho_{\star,\rm obs})$ with $b$ and $\rho_{\star,\rm true}$ held fixed), the model matches $\delta$ and $\rho_{\star,\rm obs}$ well, yet the resulting $T_{14}$ distribution is shifted by $\Delta T_{14}\approx0.35\,$h relative to the photometric-based posterior ($W_1=0.347\,\rm h$, $E=0.639\,\rm h$), which is about $6\%$ of the total transit duration. This discrepancy indicates that a likelihood based on only two observables cannot uniquely constrain the ring geometry, because multiple distinct configurations can reproduce $(\delta,\rho_{\star,\rm obs})$ while yielding different total transit durations. %RF

In our preferred configuration (the $T_{14}$-inclusive likelihood with free $\rho_{\star,\mathrm{true}}$ and $b$), the predicted distributions for all five observables agree closely with the photometric-derived posteriors: the Wasserstein distances are $W_1\lesssim4\,\rm ppm$ for $\delta$, $W_1\lesssim0.05\,\rm h$ for both $T_{14}$ and $T_{23}$, and $W_1\lesssim0.05\,\rm g\,cm^{-3}$ for $\rho_{\star,\rm obs}$. This level of agreement across the entire set of observables—including the out-of-likelihood parameter $b_{\rm obs}$—offers the strongest available evidence that these solutions reproduce ring configurations consistent with the Kepler light curves. However, the resulting posteriors for the nuisance parameters still disagree with their most stringent external prior constraints, with $\rho_{\star,\rm true}$ biased high relative to the isochrone value and $b$ biased high relative to the JWST measurement. As outlined in \autoref{sec:ring_nuisance}, these offsets are modeling artifacts caused by the tension between the orbital scale and the ring contact geometry, further exacerbated by the coarse temporal sampling of Kepler’s long-cadence data. Consequently, we are unable to isolate a fully self-consistent ring solution, which is required for a conclusive test of the ring hypothesis for the Kepler-51 system. %RF

\section{DISCUSSION AND CONCLUSIONS}
\label{sec:discussion}

The framework developed in this work was designed as an  exploratory tool to assess whether ring configurations exist that are consistent with the available Kepler-51 transit data, limited also by the $\sim 30$\,min instrumental cadence. Our Bayesian  retrievals recover ring configurations capable of reproducing the full set of \photofit-derived transit observables for both planets. However, these solutions require both $\rho_{\star,\mathrm{true}}$ and $b$ inconsistent with their current best external constraints from \citet{Berger2023} and \citet{Masuda2024}, preventing a fully self-consistent ring solution from \textit{Kepler} photometry alone. This reflects the intrinsic degeneracy of the problem: numerical convergence toward compatible ring geometries does not guarantee physical viability, which ultimately depends on tighter independent constraints. It is important to emphasize that our analysis does not constitute a direct fit of a ringed-planet transit model to the \textit{Kepler} light curves. Instead, it is based on the global transit observables set $(\observables)$, which capture only part of the information contained in the full transit morphology. A direct fit to the photometry, ideally performed jointly with a dynamical TTV model, would exploit the detailed ingress and egress shape and any residual signatures to further constrain the ring geometry and break some of the remaining parameter degeneracies in our analysis.

Although we do not obtain a fully self-consistent ring solution, our preferred retrievals for both planets converge toward compact ring systems ($f_e\approx2$--$2.7,R_p$ for \exoplanet{Kepler-51}{b} and $f_e\approx1.7$--$2.1,R_p$ for Kepler-51,d), with moderately inclined ring planes ($i_R\approx60^\circ$--$75^\circ$) and similarly moderate projected tilts ($\theta_R\approx60^\circ$--$75^\circ$). Among the explored solutions, we favor this configuration over either nearly edge-on or face-on geometries, as it is also more consistent with current theoretical expectations for ring formation and long-term stability. This interpretation further suggests a potential detectability bias. As illustrated by our forward model and the \PhotoRing\ effect map (\autoref{fig:photoring_contour}), the strongest photometric signatures arise from large, optically thick rings observed at high obliquity. The configurations preferred by our retrieval instead occupy a region of parameter space where the expected photometric signatures are comparatively weaker. If similar ring geometries are common among exoplanets, current searches would be intrinsically less sensitive to them, suggesting that at least part of the apparent lack of confirmed exoplanetary rings may arise from observational selection effects rather than a genuinely low occurrence rate.

The anomalous bulk densities under the ring hypothesis neither exclude alternative explanations for the inflated radii of the Kepler-51 planets nor are they mutually exclusive with them. Atmospheric inflation mechanisms such as high-altitude hazes, primordial H/He envelopes, or a young evolutionary state remain viable scenarios \citep{Libby-Roberts2020, Wang2019, Lammers2024}. Similarly, while the flat transmission spectra observed with HST and \textit{JWST} are commonly attributed to aerosols that mute molecular features \citep{Libby-Roberts2025}, an optically thick ring system could produce the same effect by attenuating stellar light achromatically. Distinguishing between these interpretations would require detecting wavelength-dependent ring scattering in reflected light, potentially including polarimetric signatures as predicted by \citet{Veenstra2025}. More generally, a planet hosting both an extended atmosphere and an inclined ring system would appear doubly inflated, with both effects contributing to the observed transit depth. Our ring retrievals indicate that the inferred bulk densities would increase by factors of only $\sim3$--$4$, leaving the planets substantially less dense than typical sub-Neptunes. This residual density deficit suggests that rings, if present, provide only a partial explanation for the super-puff phenomenon and that both circumplanetary material and an intrinsically low-density atmosphere are physically plausible and observationally motivated.

An important advantage of asterodensity profiling is its computational efficiency, as deriving $\rho_{\star,\mathrm{obs}}$ requires only the fundamental transit observables $(\delta, T_{14}, T_{23}, P)$. This makes the \PhotoRing\ diagnostic readily applicable to the full population of super-puff planets discovered by \textit{Kepler} and \textit{K2} without dedicated ring-transit modeling. While multi-planet architectures offer a distinct advantage by allowing the cross-testing of inter-planetary inconsistencies in \(\rho _{\star ,\mathrm{obs}}\), single-planet systems require a far more exhaustive analysis due to the lack of baseline references for contrast. Nevertheless, any system displaying such inter-planetary discrepancies, similar to those identified here, remains a prime target for detailed ring retrieval analyses. The case of \exoplanet{HIP-41378}{f} already provides the closest precedent, as its super-puff-like density was reconciled with an inclined ring system capable of inflating the apparent transit radius \citep{Akinsanmi2020}. If \PhotoRing\ offsets are found to be common among super-puffs, this would support the idea that unresolved ring systems contribute systematically to their apparent radius inflation, with important implications for sub-Neptune occurrence rates and for the theoretical models invoked to explain their formation.

More broadly, the inference framework developed here, which combines \photofit-based asterodensity profiling, a geometric ringed-planet transit model, and Bayesian retrieval, is readily transferable to other systems requiring only posterior distributions for the transit observables and an independent estimate of $\rho_{\star,\mathrm{true}}$. Its main limitation is the photometric cadence required to resolve ingress and egress morphology. In the near term, existing \textit{JWST} observations offers the most promising opportunity to test this scenario, as \textit{JWST} can directly constrain $b$ from resolved contact slopes while its spectroscopic capabilities, together with those of \textit{ARIEL}, can simultaneously probing its atmospheric properties. The upcoming high-precision observations from \textit{PLATO}$~$\citep{Matuszewski2023} for thousands of cold Jupiters will also enable systematic searches for ring signatures, improving TTV-based mass constraints, and identifying planets with anomalously low ringless densities.

A further improvement concerns the forward model itself. Its fully analytic formulation (\autoref{sec:geom_model}) was designed to enable the several likelihood evaluations required by Nested Sampling, at the expense of some physical realism. In particular, the current treatment relies on analytic expressions for the ring--stellar contact geometry. As a result, the model does not recover the ringless limit as $i_R\rightarrow90^\circ$ (edge-on), where the projected ring width vanishes and the system becomes photometrically indistinguishable from a bare planet, yet the model can still predicts non-zero biases to the transit observables. This limitation could be addressed by introducing simple analytic corrections that enforce the correct asymptotic behavior. In parallel, we are developing a fully numerical, dynamical forward model to validate and extend the analytic solution. Although such a model could also be used directly in Bayesian retrievals, its substantially higher computational cost makes the analytic formulation the preferred choice for large-scale parameter exploration.

\section*{Data Availability}

All the data necessary to reproduce the results and recreate the figures in this work—including the {\tt Jupyter} notebooks used for model training and evaluation—are publicly accessible in the {\tt GitHub} repository \url{https://github.com/seap-udea/PRisma} (PhotoRing inference and modeling algorithm). We anticipate that this project will provide a foundation for future implementations of the PR inference pipeline proposed in this paper. Contributions from the community are encouraged. %RF 

\begin{acknowledgments}

This work has been possible due to the availability of open-source, general-purpose {\tt Python} packages, including 
\  {\tt Matplotlib} \citep{Matplotlib2007}, {\tt Numpy} \citep{Numpy2020}, {\tt Scipy} \citep{virtanenSciPy10Fundamental2020}, {\tt mpi4py} \citep{dalcin2021mpi4py} and {\tt pandas} \citep{Pandas2010}, \texttt{dynesty} \citep{Speagle2020} and \texttt{emcee} \citep{ForemanMackey2013}.

\end{acknowledgments}

\newpage
\appendix

\section{Analytical Derivation of Ring Contact Positions using the Support Function}
\label{sec:appendix_support_function}

\subsection{Error Analysis of the 2015 Analytical Formula}

In the original formulation of the \PhotoRing\ effect by \citet{Zuluaga2015}, calculating the exact temporal contact points between a stellar limb and a projected exoplanetary ring requires solving a non-linear system of trigonometric equations. The previous analytical approximation assumed a highly elongated ellipse ($A \gg B$) and small contact angles, leading to geometrical artifacts at high projected tilts ($\theta_R \gtrsim 60^\circ$) and high inclinations.

To illustrate this, \autoref{fig:pr_contour_comparison} compares the \PhotoRing\ effect contours generated by the precise semi-analytical model \texttt{geotrans} with those produced by the original \texttt{exorings} analytical approximation from \citet{Zuluaga2015}. As shown, the original formulation breaks down in certain inclination ranges, creating unphysical artifacts.

\begin{figure*}[h!]
    \centering
    \includegraphics[width=0.48\textwidth]{figures/PRContour-K51b-exorings_legacy.png}
    %\includegraphics[width=0.48\textwidth]{figures/PRContour-K51b-exorings_updated.png}
    \includegraphics[width=0.48\textwidth]{figures/PRContours-K51b-geotrans.png}
    \caption{Comparison of the \PhotoRing\ effect contours for \exoplanet{Kepler-51}{b} using the original \texttt{exorings} analytical approximation from \citet{Zuluaga2015} (left) and the precise semi-analytical model \texttt{geotrans} (right). The original model exhibits geometrical artifacts at certain high projected tilts and inclinations. You can compare the panel on the right with \autoref{fig:photoring_contour}, which was obtained using the refined analytical model presented in this appendix.}
    \label{fig:pr_contour_comparison}
\end{figure*}

To provide a robust, exact-like analytical solution suitable for computationally intensive sampling (e.g., Nested Sampling), we reframe the contact geometry using the Minkowski Addition and the Support Function of a convex body \citep{schneider1993convex}.

\subsection{Theoretical Framework: The Support Function}

Geometrically, the contact between the stellar disk (a circle of radius $R_* = 1$) and the projected ring (an ellipse with semi-axes $A, B$ rotated by $\theta_R$) occurs when the distance between their centers equals the sum of their effective radii in the direction of contact. 

For a convex shape, the distance from its center to a tangent line perpendicular to a given normal vector $\vec{n}$ is defined by its Support Function, $h(\vec{n})$. For an ellipse aligned with the coordinate axes, the support function is:
\begin{equation}
    h(\vec{n}) = \sqrt{A^2 n_x^2 + B^2 n_y^2}
\end{equation}

A graphical representation of these geometric quantities and the contact positions used in this derivation is provided in \autoref{fig:analytical_model}.

\begin{figure*}[ht!]
\centering
\includegraphics[width=0.8\textwidth]{figures/PRisma-analytical-model.png}
\caption{Illustration of the quantities involved in the analytical model for two configurations: one where the rings are oriented coinciding with the coordinate system (left/top), and another where the rings are oriented with an inclination and projected tilt (right/bottom).}
\label{fig:analytical_model}
\end{figure*}

When the ellipse is rotated by an angle $\theta_R$ (the projected ring tilt), the normal vector must be rotated into the ellipse's local frame. The unit vectors defining the major and minor axes are:
\begin{equation}
    \hat{u}_A = (\cos\theta_R, \sin\theta_R), \quad \hat{u}_B = (-\sin\theta_R, \cos\theta_R) \nonumber
\end{equation}

The directional effective radius of the rotated ring is then:
\begin{equation}
    h(\vec{n}) = \sqrt{A^2 (\vec{n} \cdot \hat{u}_A)^2 + B^2 (\vec{n} \cdot \hat{u}_B)^2}
\end{equation}

\subsection{The Unperturbed Normal Approximation}

To avoid solving for the exact contact normal iteratively, we make a highly accurate physical approximation: the normal vector at the point of contact is nearly identical to the normal vector if the planet were a spherical point-mass.

For a point-mass transit with an impact parameter $b$, the contact points on the stellar limb have the coordinates $y = b$ and $x = \pm\sqrt{1-b^2}$. Therefore, the normal vectors at the egress (trailing, $R$) and ingress (leading, $L$) contacts are simply:
\begin{eqnarray}
    \vec{n}_R &=& \left( +\sqrt{1-b^2}, b \right) \\
    \vec{n}_L &=& \left( -\sqrt{1-b^2}, b \right)
\end{eqnarray}

Let $x_0 = \sqrt{1-b^2}$. We can now evaluate the Support Function of the ring in these two specific directions to obtain the effective contact radii, $h_R$ and $h_L$.

Projecting $\vec{n}_R$ onto the ellipse axes yields the egress effective radius:
\begin{eqnarray}
    \vec{n}_R \cdot \hat{u}_A &=& x_0 \cos\theta_R + b \sin\theta_R \nonumber \\
    \vec{n}_R \cdot \hat{u}_B &=& -x_0 \sin\theta_R + b \cos\theta_R \nonumber
\end{eqnarray}
Squaring and substituting into the support function gives:
\begin{eqnarray}
\label{eq:h_R}
    h_R^2 &=& A^2 (x_0 \cos\theta_R + b \sin\theta_R)^2 \nonumber \\
          & & +\, B^2 (b \cos\theta_R - x_0 \sin\theta_R)^2
\end{eqnarray}

Similarly, for the ingress effective radius ($h_L$):
\begin{eqnarray}
    \vec{n}_L \cdot \hat{u}_A &=& -x_0 \cos\theta_R + b \sin\theta_R \nonumber \\
    \vec{n}_L \cdot \hat{u}_B &=& x_0 \sin\theta_R + b \cos\theta_R \nonumber
\end{eqnarray}
Yielding:
\begin{eqnarray}
\label{eq:h_L}
    h_L^2 &=& A^2 (-x_0 \cos\theta_R + b \sin\theta_R)^2 \nonumber \\
          & & +\, B^2 (b \cos\theta_R + x_0 \sin\theta_R)^2
\end{eqnarray}

\subsection{Final Contact Equations}

With the effective geometrical radii $h_R$ and $h_L$ defined, the complex ring behaves instantaneously like a sphere of radius $h_{L,R}$. We substitute these directional radii into the standard spherical contact equations. The four contact positions along the $x$-axis are therefore given exactly by:

\noindent \textbf{Ingress Contacts (Left):}
\begin{eqnarray}
    x_{R,1} &\approx& -\sqrt{(1 + h_L)^2 - b^2} \quad \text{(First contact, external)} \\
    x_{R,2} &\approx& -\sqrt{\max(0, [1 - h_L]^2 - b^2)} \quad \text{(Second contact, internal)}
\end{eqnarray}

\noindent \textbf{Egress Contacts (Right):}
\begin{eqnarray}
    x_{R,3} &\approx& +\sqrt{\max(0, [1 - h_R]^2 - b^2)} \quad \text{(Third contact, internal)} \\
    x_{R,4} &\approx& +\sqrt{(1 + h_R)^2 - b^2} \quad \text{(Fourth contact, external)}
\end{eqnarray}

The argument of the square root for internal contacts ($x_{R,2}, x_{R,3}$) is bounded to $\ge 0$ to properly handle grazing configurations where a full interior transit does not occur.


\bibliography{references}{}
\bibliographystyle{aasjournalv7}

\end{document}